\documentclass[11pt,a4paper,oneside]{book}
\usepackage[utf8]{inputenc}
\usepackage[T1]{fontenc}
\usepackage{lmodern}
\usepackage[english]{babel}
\usepackage{amsmath,amssymb,amsthm,mathtools}
\usepackage{booktabs}
\usepackage{graphicx}
\usepackage{xcolor}
\usepackage{tikz}
\usetikzlibrary{arrows.meta,positioning,calc,shapes,shapes.geometric,decorations.pathreplacing,fit,backgrounds,trees,automata,circuits.logic.US,circuits.logic.IEC}
\usepackage{listings}
\usepackage{hyperref}
\hypersetup{colorlinks=true,linkcolor=blue!60!black,citecolor=green!50!black,urlcolor=blue!60!black}
\usepackage{enumitem}
\usepackage{geometry}
\geometry{a4paper, top=1.3cm, bottom=1.3cm, left=1.2cm, right=1.2cm, includeheadfoot}
\usepackage{setspace}
\setstretch{1.20}
\setlength{\parskip}{0.70em}
\setlength{\parindent}{1.0em}
\setlength{\abovedisplayskip}{10pt plus 2pt minus 2pt}
\setlength{\belowdisplayskip}{10pt plus 2pt minus 2pt}
\setlength{\abovedisplayshortskip}{7pt plus 2pt}
\setlength{\belowdisplayshortskip}{7pt plus 2pt}
\setlist{topsep=0.55em, partopsep=0.25em, itemsep=0.45em, parsep=0.25em}
\setlength{\textfloatsep}{10pt plus 2pt minus 2pt}
\setlength{\intextsep}{8pt plus 2pt minus 2pt}

\theoremstyle{definition}
\newtheorem{definition}{Definition}[chapter]
\newtheorem{theorem}{Theorem}[chapter]
\newtheorem{lemma}{Lemma}[chapter]
\newtheorem{corollary}{Corollary}[chapter]
\newtheorem{example}{Example}[chapter]
\newtheorem{exercise}{Exercise}[chapter]
\theoremstyle{remark}
\newtheorem{remark}{Remark}[chapter]

\lstset{basicstyle=\ttfamily\footnotesize,breaklines=true,frame=single,numbers=left,numberstyle=\tiny\color{gray},keywordstyle=\color{blue!70!black},commentstyle=\color{green!50!black},stringstyle=\color{orange!70!black},showstringspaces=false,tabsize=2, xleftmargin=1.0em, xrightmargin=0.5em, aboveskip=0.6em, belowskip=0.6em}

\title{SC1005 Digital Logic\\NTU BSc Computer Science --- Comprehensive Textbook (0--100)}
\author{NTU CS Curriculum Textbook Series\\College of Computing and Data Science}
\date{2026-08-23 --- Generated for bumbsclaw/ntu-cs-textbooks}

\begin{document}
\frontmatter
\maketitle
\tableofcontents
\mainmatter
% SC1005 --- Chapter 1: Number Systems (0->100 ground-up)
\chapter{Number Systems: From Counting to Binary and Hexadecimal}
\label{ch:numbers}

\begin{quote}
\emph{``There are 10 kinds of people: those who understand binary and those who don't.''}
Digital hardware speaks only two voltages. This chapter teaches you to count, compute, and
think fluently in that language, starting from zero---no prior digital background assumed.
\end{quote}

\noindent\fbox{\parbox{\dimexpr\linewidth-2\fboxsep-2\fboxrule}{%
\textbf{From zero: built for absolute beginners.} We assume only school arithmetic.
Every idea is picture $\to$ words $\to$ definition. By the end you will convert,
add, and negate in binary/hex as naturally as in decimal.}}
\vspace{0.4em}

\section*{Learning Outcomes}
After this chapter you will be able to:
\begin{enumerate}[label=(L\arabic*)]
    \item explain positional number systems and weight of each digit position;
    \item convert fluently among decimal, binary, octal, and hexadecimal;
    \item count binary ranges and detect overflow for $n$-bit unsigned numbers;
    \item represent signed integers via sign-magnitude, one's and two's complement;
    \item add, subtract, and reason about overflow in two's complement arithmetic.
\end{enumerate}

% ============================================================
\section{Why Not Decimal Inside a Computer?}
\label{sec:why-binary}
Decimal uses ten symbols; hardware would need ten distinct voltage levels---fragile
against noise. Two levels (\texttt{LOW} $\approx 0$\,V, \texttt{HIGH} $\approx V_{DD}$)
are maximally noise-immune. One binary digit (bit) stores one yes/no choice. Eight bits make a byte,
the standard storage atom. This is why all modern computers are binary at the bottom.

\paragraph{Positional intuition.} In decimal $472 = 4\!\times\!10^2+7\!\times\!10^1+2\!\times\!10^0$.
Replace base $10$ by $2$: each position weighs twice the one to its right. Same rules, smaller base.
The rightmost position is the least-significant bit (LSB), the leftmost the most-significant bit (MSB).

\begin{definition}[Positional representation]
In base (radix) $r$, digits $d_i\in\{0,\dots,r-1\}$ and
\[ N = d_{k-1}r^{k-1}+\cdots+d_1r^1+d_0r^0+\;.\;d_{-1}r^{-1}+\cdots \]
We write $(d_{k-1}\dots d_0.d_{-1}\dots)_r$. Decimal ($r=10$), binary ($r=2$),
octal ($r=8$), hexadecimal ($r=16$, digits $0$--$9,A$--$F$) are the four we use.
\end{definition}

\begin{figure}[h]
\centering
\begin{tikzpicture}[scale=0.82,every node/.style={font=\scriptsize}]
\foreach \i/\w/\col in {0/1/blue!12,1/2/green!14,2/4/orange!16,3/8/violet!14,4/16/red!12,5/32/teal!14,6/64/yellow!18,7/128/blue!10} {
  \node[draw,thick,rounded corners=2pt,fill=\col,minimum width=1.05cm,minimum height=0.9cm,align=center] at (\i*1.18,0.6) {$2^{\i}$\\$\w$};
  \node[font=\tiny,blue!60!black] at (\i*1.18,-0.15) {weight};
}
\node[font=\footnotesize\bfseries] at (3.6,1.7) {Binary positional weights (8-bit byte)};
\draw[->,thick,blue!60!black] (7.8,0.6)--(8.4,0.6) node[right,font=\tiny] {MSB$\to$LSB growth $\times2$};
\node[font=\scriptsize,align=center] at (3.6,-1.0) {Example: $(10110101)_2 = 128+32+16+4+1 = 181_{10}$};
\foreach \i/\b in {0/1,1/0,2/1,3/0,4/1,5/1,6/0,7/1} {
  \node[font=\footnotesize\bfseries,fill=white,inner sep=1pt] at ({(7-\i)*1.18},0.6) {\b};
}
\end{tikzpicture}
\caption{Each binary position weighs a power of two. The rightmost bit (LSB) weighs $1$, next $2$, then $4,8,\dots$.}
\label{fig:binary-weights}
\end{figure}

\section{Conversions}
\label{sec:conversions}
\paragraph{Binary $\to$ decimal:} sum weights of $1$-bits. Example: $(1101)_2 = 8+4+1=13$.

\paragraph{Decimal $\to$ binary:} repeated division by $2$, read remainders bottom-up.
$13\div2=6$ r$1$, $6\div2=3$ r$0$, $3\div2=1$ r$1$, $1\div2=0$ r$1$ $\Rightarrow 1101$.
For fractions, repeated multiply by $2$ and collect integer parts.

\paragraph{Hexadecimal shortcut.} Since $16=2^4$, group bits in fours:
$(1011\,0110)_2 = (B6)_{16}$, and $(3F)_{16}= (0011\,1111)_2$. This is why hex is ubiquitous
in datasheets, addresses, and machine code listings---compact alias for binary.

\begin{figure}[h]
\centering
\begin{tikzpicture}[scale=0.78,every node/.style={font=\scriptsize}]
\node[font=\footnotesize\bfseries] at (2.0,1.6) {Binary $\to$ Hex: group by 4 from radix point};
\foreach \i/\b in {0/1,1/0,2/1,3/1,4/0,5/1,6/1,7/0} \node[draw,thick,fill=blue!10,minimum width=0.55cm,minimum height=0.55cm] at (\i*0.62,0.7) {\b};
\draw[thick,red!70!black,rounded corners=3pt] (-0.32,0.32) rectangle (1.18,1.08);
\draw[thick,red!70!black,rounded corners=3pt] (1.24,0.32) rectangle (2.74,1.08);
\node[font=\tiny,red!70!black] at (0.43,0.0) {$1011=B$}; \node[font=\tiny,red!70!black] at (1.99,0.0) {$0110=6$};
\node at (2.7,-0.7) {$\Rightarrow (B6)_{16}$};
\node[font=\footnotesize\bfseries] at (7.0,1.6) {Hex $\to$ Binary: expand each digit};
\foreach \i/\h in {0/3,1/F} { \node[draw,thick,fill=orange!18,minimum width=0.7cm,minimum height=0.55cm] at (6.2+\i*1.1,0.7) {\h}; }
\node[font=\tiny] at (6.2,0.0) {$0011$}; \node[font=\tiny] at (7.3,0.0) {$1111$};
\draw[->,thick,blue!60!black] (5.0,0.7)--(5.7,0.7);
\draw[->,thick,blue!60!black] (7.8,0.7)--(8.3,0.7) node[right] {$(00111111)_2$};
\draw[->,thick,teal!70!black] (2.0,-1.1)--(2.0,-1.45) node[below,font=\tiny] {$\div2$ method $\leftrightarrow$ $\times2$ for fractions};
\end{tikzpicture}
\caption{Binary--hex conversion is lossless grouping by four bits; decimal conversion uses repeated division/multiplication.}
\label{fig:hex-grouping}
\end{figure}

\begin{example}[Conversions]
Convert $(2A.4)_{16}$ to decimal: $2\!\times\!16+10+4/16=42.25$.
Convert $42.25$ to binary: $42=(101010)_2$, $0.25=(.01)_2$, so $101010.01_2$---check hex $2A.4$ matches.
Octal groups by $3$ bits: $(101101)_2=(55)_8$ since $16=2^4$, $8=2^3$.
\end{example}

\section{Unsigned Range and Overflow}
An $n$-bit unsigned integer spans $0$ to $2^n-1$. With $n=8$, $0$--$255$; $n=32$, $0$--$4\,294\,967\,295$.
Adding $1$ to $111\ldots1$ wraps to $00\ldots0$ (overflow / carry out). Hardware signals this with a carry flag.

\begin{figure}[h]
\centering
\begin{tikzpicture}[scale=0.82,every node/.style={font=\scriptsize}]
\draw[thick,fill=blue!4] (0,0) circle (1.45);
\foreach \i/\b/\d in {0/000/0,1/001/1,2/010/2,3/011/3,4/100/4,5/101/5,6/110/6,7/111/7} {
  \pgfmathsetmacro{\a}{90-\i*45}
  \node[draw,circle,fill=orange!\the\numexpr 10+\i*8\relax,inner sep=2pt,font=\tiny] at (\a:1.0) {\b};
  \node[font=\tiny] at (\a:1.75) {\d};
}
\draw[->,thick,red!70!black] (60:1.25) arc (60:20:1.25) node[midway,right,font=\tiny] {+1 wraps};
\node[font=\footnotesize\bfseries] at (0,2.25) {3-bit unsigned wheel: $0$--$7$ then wrap};
\node[font=\tiny,align=center] at (0,-2.25) {Adding 1 moves clockwise;\\ overflow = carry beyond $n$ bits};
\begin{scope}[xshift=5.2cm]
\draw[thick,fill=green!12,rounded corners=2pt] (0,0.4) rectangle (3.2,-0.4);
\foreach \x in {0,0.8,1.6,2.4,3.2} \draw[thick] (\x,0.4)--(\x,-0.4);
\node[font=\tiny] at (0.4,0) {$0$}; \node[font=\tiny] at (1.2,0) {$1$}; \node[font=\tiny] at (2.0,0) {$\cdots$}; \node[font=\tiny] at (2.8,0) {$2^n{-}1$};
\node[font=\tiny,align=center] at (1.6,-0.85) {$n$ bits $\to$ $2^n$ values};
\node[font=\footnotesize\bfseries] at (1.6,1.05) {Unsigned range};
\end{scope}
\end{tikzpicture}
\caption{Unsigned $n$-bit numbers live on a wrap-around wheel. The extra carry beyond $n$ bits is the overflow flag.}
\label{fig:number-wheel}
\end{figure}

\section{Signed Numbers and Two's Complement}
Three historical schemes for signed $n$-bit integers:
\begin{itemize}
    \item \textbf{Sign-magnitude:} MSB is sign ($0=+,1=-$), rest is magnitude. Two zeros! Range $-(2^{n-1}-1)$ to $+(2^{n-1}-1)$.
    \item \textbf{One's complement:} Negate by flipping all bits. Also two zeros. Rare today.
    \item \textbf{Two's complement} (universal): Negate by invert+1. One zero, range $-2^{n-1}$ to $2^{n-1}-1$. Addition is same as unsigned---hardware's reason for winning.
\end{itemize}

\begin{definition}[Two's complement]
For $n$ bits, value of bit-vector $b_{n-1}\dots b_0$ is
\[ V = -b_{n-1}2^{\,n-1}+\sum_{i=0}^{n-2} b_i2^i. \]
Negation: $\; -N = \overline{N}+1$ (bitwise NOT plus one). Subtracting $A-B$ is $A+\overline{B}+1$.
\end{definition}

\begin{example}[Two's complement arithmetic]
$n=4$: $+3=0011$, $-3=1101$ (invert $0011\to1100$+1). $3+(-3)=0011+1101=0000$ (carry discarded) correct.
$5+6=0101+0110=1011$ which is $-5$ after wrap---signed overflow since two positives yielded negative.
$4$-bit range is $-8$ to $+7$; $1011$ as signed is $-5$, as unsigned is $11$.
\end{example}

\begin{theorem}[Overflow rule]
Signed overflow in two's complement addition occurs iff carry into sign bit $\neq$ carry out of sign bit,
equivalently: adding two same-sign numbers yields opposite sign. Unsigned overflow is simply carry out $=1$.
\end{theorem}

\section{Binary Codes and Arithmetic Preview}
Not all bit patterns are numbers. \textbf{BCD} encodes each decimal digit in $4$ bits (wasteful but decimal-exact, used in calculators).
\textbf{Gray code} changes one bit per step---ideal for rotary encoders and Karnaugh maps (no glitch).
\textbf{ASCII/Unicode} map bit patterns to characters; $(41)_{16}$ = `A', $(30)_{16}$ = `0'.
Binary addition is column-wise with carries, exactly like decimal but base $2$: $1+1=0$ carry $1$, $1+1+1=1$ carry $1$.


\section{Fractions, Floating-Point Preview, and Binary Arithmetic}
Binary fractions mirror decimal: $(0.101)_2 = 1/2+1/8=0.625$. Not every decimal fraction terminates in binary
($0.1_{10}=0.00011\overline{0011}_2$ recurring), which is why floating-point rounding matters---preview for
computer arithmetic courses. Binary addition is column-wise:
\begin{center}
\begin{tabular}{rccccc}
 & 1 & 0 & 1 & 1 \\ % 11
+& 0 & 1 & 1 & 0 \\ % 6
\hline
carry &1&1&0&0\\
result&1&0&0&0&1 % 17 - need 5 bits? Actually 11+6=17=10001
\end{tabular}
\end{center}
Carry propagates left; the extra carry beyond $n$ bits is exactly the unsigned overflow flag $C$.
Subtraction via two's complement: $A-B=A+\bar B+1$, so one adder does both operations with an invert-and-carry-in control.

\begin{example}[Overflow table]
$n=4$ signed: $7+1=0111+0001=1000$ ($-8$) overflow; $-8-1=1000+1111=0111$ ($+7$) overflow.
Unsigned, $15+1=1111+0001=0000$ carry $1$---correct modulo $2^4$ wrap.
\end{example}

% ============================================================

% ============================================================
\section{Chapter Notes}
Positional notation dates to Hindu-Arabic mathematics; binary championed by Leibniz (1703).
Two's complement emerged with early stored-program machines (EDSAC, 1949) for its single-adder economy.
Hex became standard with the IBM System/360 (1964) to print 32-bit words compactly.

% ============================================================
\section{Fixed-Point vs.\ Floating Preview}
Fixed-point binary $(b_{k-1}\dots b_0.b_{-1}\dots b_{-m})_2$ has uniform step $2^{-m}$ and range $\approx2^k$.
More bits after the point = finer resolution; more before = larger range. For large dynamic range,
floating-point (sign $\times$ mantissa $\times2^{\text{exponent}}$) trades uniform precision for wide range---full
treatment in computer arithmetic courses. Key takeaway: not all reals are representable; rounding error is inevitable
and must be bounded (e.g.\ unit in last place, ULP).

% ============================================================

\section{Exercises}
\label{sec:ch01-exercises}
\begin{enumerate}[label=E1.\arabic*, leftmargin=*, itemsep=0.8em]
    \item \textbf{Conversions.} Convert $(101101.101)_2$ to decimal, hex, and octal. Convert $(3A7)_{16}$ and $(275)_8$ to binary. Show steps (division/multiplication).
    \item \textbf{Two's complement.} For $n=8$, give binary for $+73$, $-73$, $-128$, $+127$. Compute $73+(-98)$ in binary and flag signed and unsigned overflow.
    \item \textbf{Range and wrap.} An $n=5$ unsigned counter increments by $1$ from $0$. After how many increments does it first read $0$ again? Repeat for signed two's complement starting at $+15$ and incrementing.
    \item \textbf{Hex dump reading.} A memory dump shows \texttt{DE AD BE EF}. Write as binary, as four unsigned bytes decimal, and as two 16-bit big-endian unsigned values. What about little-endian?
    \item \textbf{Gray vs.\ binary.} Construct 3-bit Gray code, show adjacent words differ by one bit, and explain why a 2-bit transition $011\to100$ in binary could briefly read $111$ on skewed wires.
\end{enumerate}
% End of Chapter 1

% SC1005 --- Chapter 2: Boolean Algebra & Karnaugh Maps
\chapter{Boolean Algebra and Karnaugh Maps}
\label{ch:boolean}

\begin{quote}
\emph{``George Boole made logic into algebra; Shannon made algebra into circuits.''}
This chapter turns truth tables into equations, equations into minimal circuits, and
minimal circuits into patterns you can see on a map---no prior logic assumed.
\end{quote}

\noindent\fbox{\parbox{\dimexpr\linewidth-2\fboxsep-2\fboxrule}{%
\textbf{From zero: logic from scratch.} We define propositions, connectives, and axioms
before any algebra. Each law gets a Venn or gate picture first, then the identity.}}
\vspace{0.4em}

\section*{Learning Outcomes}
After this chapter you will be able to:
\begin{enumerate}[label=(L\arabic*)]
    \item evaluate and manipulate Boolean expressions using axioms and theorems;
    \item obtain canonical SOP/POS forms from truth tables;
    \item simplify functions algebraically and with Karnaugh maps up to 4 variables;
    \item identify prime implicants, essential primes, and don't-care handling;
    \item describe NAND/NOR universality and two-level cost metrics.
\end{enumerate}

% ============================================================
\section{Boolean Algebra: Axioms and Intuition}
\label{sec:boolean-axioms}
A \emph{Boolean variable} takes values in $\{0,1\}$ (also $\{\text{false},\text{true}\}$).
Operations: $\cdot$ (AND), $+$ (OR), $'$ or $\overline{\phantom{x}}$ or $\neg$ (NOT).
Think sets: AND=intersection, OR=union, NOT=complement---the Venn picture \emph{is} the law.

\begin{definition}[Huntington axioms]
For all $x,y,z\in\{0,1\}$ with constants $0,1$:
$ x+0=x,\;x\cdot1=x$ (identity); $x+x'=1,\;x\cdot x'=0$ (complement);
$x+y=y+x,\;x\cdot y=y\cdot x$ (commutative); $x+(y\cdot z)=(x+y)\cdot(x+z)$ and
$x\cdot(y+z)=x\cdot y+x\cdot z$ (distributive); plus closure.
Duality: swap $+\leftrightarrow\cdot$ and $0\leftrightarrow1$ to get a dual theorem for free.
\end{definition}

\begin{theorem}[Basic theorems]
Idempotent $x+x=x$, $x\cdot x=x$; null $x+1=1$, $x\cdot0=0$; involution $(x')'=x$;
absorption $x+x\cdot y=x$, $x\cdot(x+y)=x$; De Morgan $(x+y)'=x'\cdot y'$, $(x\cdot y)'=x'+y'$;
consensus $x\cdot y+x'\cdot z+y\cdot z=x\cdot y+x'\cdot z$.
\end{theorem}

\begin{figure}[h]
\centering
\begin{tikzpicture}[scale=0.78,every node/.style={font=\scriptsize}]
% Venn AND / OR / NOT
\begin{scope}
\draw[thick,fill=blue!10] (0,0) circle (0.9);
\draw[thick,fill=orange!18] (0.7,0) circle (0.9);
\node at (-0.55,-1.25) {$x$}; \node at (1.25,-1.25) {$y$};
\begin{scope}
\clip (0,0) circle (0.9); \fill[green!40] (0.7,0) circle (0.9);
\end{scope}
\draw[thick] (0,0) circle (0.9); \draw[thick] (0.7,0) circle (0.9);
\node[font=\footnotesize\bfseries] at (0.35,1.45) {$x\cdot y$ (AND)};
\node[font=\tiny] at (0.35,0) {overlap};
\end{scope}
\begin{scope}[xshift=3.2cm]
\draw[thick,fill=blue!10] (0,0) circle (0.9);
\draw[thick,fill=orange!18] (0.7,0) circle (0.9);
\fill[green!22,even odd rule] (0,0) circle (0.9) (0.7,0) circle (0.9);
\begin{scope}
\clip (0,0) circle (0.9); \fill[green!40] (0.7,0) circle (0.9);
\end{scope}
\draw[thick] (0,0) circle (0.9); \draw[thick] (0.7,0) circle (0.9);
\node[font=\footnotesize\bfseries] at (0.35,1.45) {$x+y$ (OR)};
\end{scope}
\begin{scope}[xshift=6.4cm]
\draw[thick,fill=gray!12] (-0.1,-0.1) rectangle (1.6,1.1);
\draw[thick,fill=blue!14] (0.5,0.2) circle (0.62);
\fill[white] (0.5,0.2) circle (0.62); \fill[green!30] (-0.1,-0.1) rectangle (1.6,1.1);
\begin{scope}\clip (-0.1,-0.1) rectangle (1.6,1.1); \fill[white] (0.5,0.2) circle (0.62);\end{scope}
\draw[thick] (0.5,0.2) circle (0.62); \draw[thick] (-0.1,-0.1) rectangle (1.6,1.1);
\node[font=\footnotesize\bfseries] at (0.7,1.45) {$x'$ (NOT)};
\node[font=\tiny] at (0.7,-0.4) {outside $x$};
\end{scope}
% De Morgan arrow
\draw[->,thick,red!70!black] (2.1,-1.7)--(2.8,-1.7) node[midway,below,font=\tiny] {De Morgan swaps overlap $\leftrightarrow$ outside};
\end{tikzpicture}
\caption{Venn intuition: AND is overlap, OR is union, NOT is outside. De Morgan's laws are visually obvious.}
\label{fig:venn-boolean}
\end{figure}

\section{Canonical Forms: SOP and POS}
\label{sec:canonical}
A \emph{minterm} $m_i$ is an AND of all literals (each variable complemented or not) that is $1$ on exactly one row.
A \emph{maxterm} $M_i$ is an OR that is $0$ on exactly one row. Example for $x,y,z$:
$m_5=\bar x y\bar? $ no---with order $x\,y\,z$, $m_5=x'\cdot y\cdot z'$? Check: binary $101=5$.

Any function has two canonical forms:
\[ F=\sum m(\text{rows where }F=1)\quad\text{(SOP)},\qquad
   F=\prod M(\text{rows where }F=0)\quad\text{(POS)}. \]
We use $\sum m(1,3,6)$ as shorthand.

\begin{figure}[h]
\centering
\begin{tikzpicture}[scale=0.78,every node/.style={font=\scriptsize}]
% Gate symbols: AND, OR, NOT, NAND
\node[font=\footnotesize\bfseries] at (3.8,2.0) {Basic gates (color = function family)};
% AND
\begin{scope}[xshift=0cm]
\draw[thick,fill=blue!14] (0,0) -- (0,1) -- (0.6,1) arc (90:-90:0.5) -- (0.6,0) -- cycle;
\draw[thick] (-0.5,0.75)--(0,0.75) node[left,font=\tiny] {$x$};
\draw[thick] (-0.5,0.25)--(0,0.25) node[left,font=\tiny] {$y$};
\draw[thick] (1.1,0.5)--(1.6,0.5) node[right] {$x\cdot y$};
\node[font=\tiny,align=center] at (0.55,-0.4) {AND};
\end{scope}
% OR
\begin{scope}[xshift=2.5cm]
\draw[thick,fill=orange!18] (0,0) .. controls (0.2,0.5) .. (0,1) -- (0.35,1) .. controls (0.7,0.5) .. (0.35,0) -- cycle;
\draw[thick] (-0.5,0.75)--(0,0.75); \draw[thick] (-0.5,0.25)--(0,0.25);
\draw[thick] (0.75,0.5)--(1.25,0.5) node[right] {$x+y$};
\node[font=\tiny] at (0.35,-0.4) {OR};
\end{scope}
% NOT
\begin{scope}[xshift=4.8cm]
\draw[thick,fill=green!16] (0,0) -- (0,1) -- (0.7,0.5) -- cycle;
\draw[thick,fill=white] (0.7,0.5) circle (0.11);
\draw[thick] (-0.5,0.5)--(0,0.5) node[left,font=\tiny] {$x$};
\draw[thick] (0.81,0.5)--(1.3,0.5) node[right] {$x'$};
\node[font=\tiny] at (0.35,-0.4) {NOT};
\end{scope}
% NAND
\begin{scope}[xshift=6.8cm]
\draw[thick,fill=violet!14] (0,0) -- (0,1) -- (0.6,1) arc (90:-90:0.5) -- (0.6,0) -- cycle;
\draw[thick,fill=white] (1.1,0.5) circle (0.11);
\draw[thick] (-0.5,0.75)--(0,0.75); \draw[thick] (-0.5,0.25)--(0,0.25);
\draw[thick] (1.21,0.5)--(1.65,0.5) node[right] {$(xy)'$};
\node[font=\tiny] at (0.55,-0.4) {NAND};
\end{scope}
% Legend
\node[font=\tiny,fill=yellow!12,rounded corners=2pt,inner sep=2pt] at (3.8,-1.1) {Bubbled output = complemented; NAND/NOR are universal (any function from one gate type)};
\end{tikzpicture}
\caption{Gate symbols with color families. The small bubble denotes inversion; NAND follows AND with a bubble.}
\label{fig:gates}
\end{figure}

\section{Karnaugh Maps}
\label{sec:kmaps}
Algebraic minimization is error-prone. A \emph{Karnaugh map} places minterms so adjacent cells differ by one literal (Gray order). Grouping $2^k$ adjacent $1$s eliminates $k$ literals.

Rules: groups must be rectangular, size power of two, wrap around edges, cover all $1$s, make as large and as few as possible. A \emph{prime implicant} cannot be enlarged; an \emph{essential} one covers a $1$ no other prime covers.

\begin{figure}[h]
\centering
\begin{tikzpicture}[scale=0.78,every node/.style={font=\scriptsize}]
% 3-var K-map example F = sum m(0,1,4,5,6)
\node[font=\footnotesize\bfseries] at (2.2,3.0) {3-var K-map: $F(x,y,z)=\sum m(0,1,4,5,6)$};
% grid 2x4
\foreach \col in {0,...,3} \foreach \row in {0,1} \draw[thick] (\col*1.15,\row*1.05) rectangle (\col*1.15+1.15,\row*1.05+1.05);
% headers yz Gray: 00 01 11 10
\node at (0.575,2.55) {$00$}; \node at (1.725,2.55) {$01$}; \node at (2.875,2.55) {$11$}; \node at (4.025,2.55) {$10$};
\node at (4.025,2.22) {$yz$}; \node at (-0.45,1.575) {$x$}; \node at (-0.45,1.05) {$0$}; \node at (-0.45,0.0) {$1$};
% fill minterms
\node at (0.575,1.575) {1}; \node at (1.725,1.575) {1}; \node at (2.875,1.575) {0}; \node at (4.025,1.575) {0};
\node at (0.575,0.525) {1}; \node at (1.725,0.525) {1}; \node at (2.875,0.525) {0}; \node at (4.025,0.525) {1};
% groups: top row pair (0,1) blue, bottom row pair (4,5) green, vertical pair (4,6) orange wrapping concept (col0 col3 not contiguous due to layout but 4-6 needs explanation) actually 4 is col0 row1, 6 is col3 row1 -> wrap group
\draw[thick,blue!60!black,rounded corners=3pt] (0.05,1.05) rectangle (2.30,2.12);
\draw[thick,green!60!black,rounded corners=3pt] (0.05,0.0) rectangle (2.30,1.05);
\draw[thick,orange!70!black,rounded corners=3pt] (-0.05,-0.05) rectangle (1.20,1.10); % placeholder for 0,4
% Better: group col0 both rows (0,4) orange, group 0,1,4,5 quad blue would be better but keep simple
% Redraw correctly: quad 0,1,4,5
\draw[thick,red!70!black,rounded corners=4pt] (0.0,0.0) rectangle (2.32,2.12);
\node[font=\tiny,red!70!black] at (1.15,-0.35) {quad $y'=0$ $\to$ $y'$};
% Single extra: 6 alone vs with 4
\draw[thick,violet!70!black,rounded corners=3pt,dashed] (0.0,0.0) rectangle (1.20,1.10);
% Actually show group (4,6) wrapping? Use arrow
\node[font=\tiny,align=left] at (6.2,1.7) {Groups:};
\node[font=\tiny,align=left] at (6.2,1.15) {red quad $0,1,4,5$\\$=y'$};
\node[font=\tiny,align=left] at (6.2,0.45) {orange pair $4,6$\\$=x\,z'$};
\node[font=\tiny] at (6.2,-0.15) {$\;F=y'+xz'$};
% Don't-care example mini
\end{tikzpicture}
\caption{K-map with Gray-ordered columns. The red quad of four eliminates two literals to $y'$; the pair $4$--$6$ gives $xz'$. Result $F=y'+xz'$. Adjacent includes wrap-around.}
\label{fig:kmap}
\end{figure}

\paragraph{Don't-cares.} Input combinations that never occur are marked $\times$ and may be taken as $0$ or $1$ whichever yields larger groups---free simplifiers.

\paragraph{Two-level cost.} After minimization, cost $\approx$ number of gates + literals. Fewer product terms and fewer literals per term wins for SOP (and dually for POS).

% ============================================================
\section{Chapter Notes}
Boole's \emph{Laws of Thought} (1854) and Shannon's 1937 master's thesis linking switching circuits to Boolean algebra. Karnaugh (1953) and Veitch diagrams systematized minimization before Quine--McCluskey automation.

% ============================================================
\section{Exercises}
\label{sec:ch02-exercises}
\begin{enumerate}[label=E2.\arabic*, leftmargin=*, itemsep=0.8em]
    \item \textbf{Algebra.} Simplify $F=x y + x' z + y z$ to $x y + x' z$ using consensus, then verify by truth table.
    \item \textbf{Canonical forms.} For $F(a,b,c)=\sum m(0,2,3,5,7)$, write canonical POS and simplify via K-map to minimal SOP. List prime implicants and essential primes.
    \item \textbf{Don't-cares.} $F(w,x,y,z)=\sum m(1,3,7,11,15)+\sum d(0,2,5)$. Minimize with don't-cares; show how treating each $\times$ as $1$ or $0$ changes the cover.
    \item \textbf{NAND universality.} Using only 2-input NAND gates, realize $F=ab+cd$. Draw the gate diagram and count gates.
    \item \textbf{Duality check.} State the dual of absorption $x+x y=x$ and prove it. Use duality to obtain the POS form from an SOP minimization without remapping.
\end{enumerate}
% End of Chapter 2

% SC1005 --- Chapter 3: Combinational Logic (0->100)
\chapter{Combinational Logic: Adders, Multiplexers, Decoders}
\label{ch:combinational}

\begin{quote}
\emph{``Combinational means memoryless---output is a pure function of present inputs.''}
From gates we build arithmetic, selection, and routing---the data-path building blocks
every processor uses. No sequential knowledge assumed; we start from truth tables.
\end{quote}

\noindent\fbox{\parbox{\dimexpr\linewidth-2\fboxsep-2\fboxrule}{%
\textbf{From zero: gates $\to$ blocks.} We define combinational vs.\ sequential,
then derive each block from its truth table $\to$ Boolean expression $\to$ circuit.}}
\vspace{0.4em}

\section*{Learning Outcomes}
After this chapter you will be able to:
\begin{enumerate}[label=(L\arabic*)]
    \item distinguish combinational from sequential circuits;
    \item design half/full adders and ripple-carry adders from truth tables;
    \item use multiplexers as logic and routing elements;
    \item decode addresses with decoders and encode with priority encoders;
    \item analyse hazards and two-level vs.\ multi-level cost.
\end{enumerate}

% ============================================================
\section{Combinational Building Blocks}
\label{sec:comb-def}
A \emph{combinational} circuit has no memory: output at time $t$ depends only on inputs at $t$.
No feedback, no clock. Given $n$ inputs and $m$ outputs it realises a Boolean function
$f:\{0,1\}^n\to\{0,1\}^m$, fully specified by a truth table.

Design flow: (1) truth table, (2) SOP/POS or K-map to minimise (Ch.~\ref{ch:boolean}),
(3) map to available gates (NAND/NOR universality), (4) check delay and hazards.

\section{Adders: Half, Full, Ripple-Carry}
\label{sec:adders}
Adding is the first non-trivial combinational function. One column of binary addition
produces a sum bit and a carry.

\begin{definition}[Half and full adder]
\emph{Half adder} (HA): inputs $a,b$; outputs $S=a\oplus b$, $C_{\text{out}}=a\cdot b$
(adds two bits, no carry-in). \emph{Full adder} (FA): inputs $a,b,C_{\text{in}}$;
$S=a\oplus b\oplus C_{\text{in}}$, $C_{\text{out}}=ab+C_{\text{in}}(a\oplus b)$.
\end{definition}

\begin{figure}[h]
\centering
\begin{tikzpicture}[scale=0.82,every node/.style={font=\scriptsize}]
% Half adder
\node[font=\footnotesize\bfseries] at (1.6,2.0) {Half adder};
\draw[thick,fill=blue!12,rounded corners=2pt] (0.4,0.1) rectangle (1.9,1.35);
\node at (1.15,1.05) {XOR}; \node at (1.15,0.45) {AND};
\draw[->,thick] (0,1.2)--(0.4,1.2) node[left] {$a$};
\draw[->,thick] (0,0.4)--(0.4,0.4) node[left] {$b$};
\draw[->,thick] (1.9,1.05)--(2.8,1.05) node[right] {$S$};
\draw[->,thick] (1.9,0.45)--(2.8,0.45) node[right] {$C_o$};
\node[font=\tiny,fill=yellow!12,rounded corners=1pt] at (1.15,-0.35) {$S=a\oplus b,\;C_o=ab$};
% Full adder
\begin{scope}[xshift=4.2cm]
\node[font=\footnotesize\bfseries] at (1.6,2.0) {Full adder (2 HAs + OR)};
\draw[thick,fill=green!12,rounded corners=2pt] (0.2,0.1) rectangle (1.35,1.5);
\node[font=\tiny] at (0.78,1.2) {HA1};
\draw[thick,fill=orange!14,rounded corners=2pt] (1.65,0.1) rectangle (2.8,1.0);
\node[font=\tiny] at (2.22,0.7) {HA2};
\draw[->,thick] (-0.3,1.35)--(0.2,1.35) node[left] {$a$};
\draw[->,thick] (-0.3,0.85)--(0.2,0.85) node[left] {$b$};
\draw[->,thick] (1.8,1.45)--(2.6,1.45) node[above,font=\tiny] {$C_{\text{in}}$};
\draw[->,thick] (1.35,0.85)--(1.65,0.85);
\draw[->,thick] (1.35,0.35)--(1.65,0.35);
\draw[->,thick] (2.8,0.7)--(3.4,0.7) node[right] {$S$};
% OR for carry
\draw[thick,fill=violet!12] (1.0,-0.55) rectangle (2.0,-0.05);
\node[font=\tiny] at (1.5,-0.30) {OR};
\draw[->,thick] (0.7,0.2)--(1.0,-0.15); \draw[->,thick] (2.0,0.2)--(2.0,-0.15);
\draw[->,thick] (2.0,-0.30)--(3.4,-0.30) node[right] {$C_{\text{out}}$};
\end{scope}
\end{tikzpicture}
\caption{Left: half adder. Right: full adder built from two half adders plus OR for the carry. Coloured blocks match function.}
\label{fig:half-full-adder}
\end{figure}

Chaining FAs gives a \emph{ripple-carry adder} (RCA): $C_{\text{out}}$ of bit $i$ feeds
$C_{\text{in}}$ of bit $i+1$. Simple but carry ripples---delay $O(n)$ gate levels.
Faster adders (carry-lookahead, Ch.~\ref{ch:arithmetic}) compute carries in parallel.

\begin{figure}[h]
\centering
\begin{tikzpicture}[scale=0.78,every node/.style={font=\scriptsize}]
\node[font=\footnotesize\bfseries] at (4.0,1.85) {4-bit ripple-carry adder};
\foreach \i in {0,1,2,3} {
  \pgfmathsetmacro{\x}{\i*1.95}
  \node[draw,thick,fill=blue!\the\numexpr 10+\i*8\relax,minimum width=1.55cm,minimum height=0.95cm,rounded corners=2pt] at (\x+0.9,0.6) {FA$_{\i}$};
  \node[font=\tiny] at (\x+0.9,0.9) {$a_{\i}\;b_{\i}$};
  \node[font=\tiny] at (\x+0.9,0.3) {$S_{\i}$};
  \draw[->,thick,blue!60!black] (\x+0.35,1.35)--(\x+0.55,1.05) node[above,font=\tiny] {$a_{\i}$};
  \draw[->,thick,red!60!black] (\x+1.25,1.35)--(\x+1.05,1.05) node[above,font=\tiny] {$b_{\i}$};
}
\foreach \i in {0,1,2} {
  \pgfmathsetmacro{\x}{\i*1.95}
  \draw[->,thick,orange!70!black] (\x+1.65,0.15)--(\x+1.95,0.15) node[midway,above,font=\tiny] {$C_{\i+1}$};
}
\node[font=\tiny] at (0.0,0.15) {$C_0=0$} ; \draw[->,thick] (-0.55,0.15)--(0.12,0.15);
\draw[->,thick] (7.65,0.15)--(8.15,0.15) node[right,font=\tiny] {$C_4$};
\node[font=\tiny,fill=yellow!12,rounded corners=1pt] at (4.0,-0.75) {Delay $=4\cdot t_{\text{FA}}$; carry ripples LSB $\to$ MSB};
\end{tikzpicture}
\caption{Ripple-carry adder: four full adders chained. The carry path (orange) determines critical delay.}
\label{fig:ripple-adder}
\end{figure}

\section{Multiplexers and Demultiplexers}
\label{sec:mux}
A $2^n{\times}1$ \emph{multiplexer} (MUX) selects one of $2^n$ data inputs via $n$ select lines;
a \emph{demultiplexer} does the reverse. A MUX is a universal logic block: any $n$-variable
function can be realised by wiring its truth table to data inputs and variables to selects.

\begin{figure}[h]
\centering
\begin{tikzpicture}[scale=0.82,every node/.style={font=\scriptsize}]
% 4:1 MUX
\draw[thick,fill=blue!10] (0,0) -- (0,2.2) -- (0.6,2.0) -- (0.6,0.2) -- cycle;
\foreach \y/\lbl in {1.85/$D_0$,1.35/$D_1$,0.85/$D_2$,0.35/$D_3$} \draw[->,thick] (-0.9,\y)--(0,\y) node[left,font=\tiny] {\lbl};
\draw[->,thick,red!70!black] (0.15,-0.5)--(0.15,0.05) node[below,font=\tiny] {$S_1$};
\draw[->,thick,red!70!black] (0.45,-0.5)--(0.45,0.05) node[below,font=\tiny] {$S_0$};
\draw[->,thick] (0.6,1.1)--(1.4,1.1) node[right] {$Y$};
\node[font=\tiny,align=center] at (0.3,1.1) {4:1\\MUX};
\node[font=\footnotesize\bfseries] at (0.3,2.65) {4-to-1 MUX};
% Equation
\node[font=\tiny,fill=orange!10,rounded corners=2pt] at (3.2,1.1) {$Y=S_1'S_0'D_0+S_1'S_0D_1+S_1S_0'D_2+S_1S_0D_3$};
% 2:1 MUX as gate
\begin{scope}[xshift=5.8cm]
\draw[thick,fill=green!10] (0,0) -- (0,1.6) -- (0.5,1.4) -- (0.5,0.2) -- cycle;
\draw[->,thick] (-0.7,1.25)--(0,1.25) node[left,font=\tiny] {$A$};
\draw[->,thick] (-0.7,0.35)--(0,0.35) node[left,font=\tiny] {$B$};
\draw[->,thick,red!70!black] (0.15,-0.4)--(0.15,0.05) node[below,font=\tiny] {$S$};
\draw[->,thick] (0.5,0.8)--(1.1,0.8) node[right] {$Y$};
\node[font=\tiny] at (0.25,0.8) {2:1};
\node[font=\tiny,fill=yellow!12] at (0.25,-0.85) {$Y=S'A+SB$};
\node[font=\footnotesize\bfseries] at (0.25,2.05) {2-to-1 MUX};
\end{scope}
\end{tikzpicture}
\caption{Multiplexers: $4{:}1$ selects one of four data lines; $2{:}1$ realises $Y=S'A+SB$. Select lines in red.}
\label{fig:mux}
\end{figure}

\section{Decoders, Encoders, and Tri-State}
A $n$-to-$2^n$ \emph{decoder} asserts exactly one of $2^n$ outputs for each input code
(address decoding). A \emph{priority encoder} reverses this, reporting the highest-priority active input.
\emph{Tri-state} buffers allow bus sharing: output is $0$, $1$, or high-impedance $Z$ (disconnected).

Hazards: combinational glitches when inputs change (static hazard: momentary $0$ or $1$ spike).
K-map grouping with overlapping prime implicants or adding hazard covers fixes it for two-level logic.


\section{Comparators, Code Converters, and Hierarchical Design}
Beyond adders and MUXes, two more combinational staples appear everywhere.
A \emph{magnitude comparator} reports $A>B$, $A=B$, $A<B$ via subtraction without storing the difference:
$A-B$ sign and zero flags give the answer. A \emph{code converter} (e.g.\ binary to Gray or BCD to 7-segment)
is simply a truth table mapped to SOP with don't-cares for unused codes.

Hierarchical design builds large blocks from small: an 8-bit RCA from two 4-bit RCAs, a 16-to-1 MUX from five 4-to-1 MUXes,
a 5-to-32 decoder from 2-to-4 plus 3-to-8 decoders with enables. Enables ($\overline{EN}$) are the glue---inactive blocks
output $0$ or $Z$, so only the selected block drives the bus.
\begin{example}[7-segment decoder]
Hex digit $h_3h_2h_1h_0$ drives segments $a$--$g$. Row $h=0000$ lights $a,b,c,d,e,f$ but not $g$ (displays `0').
K-maps for each segment share product terms---multi-output minimization saves gates versus per-output solo maps.
\end{example}

% ============================================================

% ============================================================
\section{Chapter Notes}
Claude Shannon's switching algebra (1938) first mapped Boolean functions to relay circuits.
MUX-based design became popular with MSI chips (74151, 74138) in the 1970s TTL era.

% ============================================================
\section{Timing and Hazards in Combinational Logic}
Gate delay matters. The \emph{critical path} is the longest input-to-output chain; its delay sets the minimum clock period
for the surrounding sequential system. A \emph{static hazard} is a momentary glitch when one input changes but the output
should stay constant---caused by unequal path delays reconverging (e.g.\ $F=x\bar x$ momentarily $1$ during transition).
A \emph{dynamic hazard} toggles multiple times before settling.

K-maps naturally reveal static hazards: any adjacent $1$s not covered by a common implicant is a potential glitch when
crossing between them. Adding the \emph{consensus} implicant that bridges the adjacency removes the hazard (at the cost of one extra term).
Modern synthesis with multi-level logic may still introduce hazards---hazard-free design or registered outputs (Ch.~\ref{ch:sequential}) solves it.

% ============================================================

\paragraph{Bus and tri-state.} Multiple combinational blocks sharing a bus use \emph{tri-state} buffers ($0/1/Z$).
Only one driver is enabled at a time; the bus is a wired-OR of tri-states. Modern CMOS prefers MUX-based buses internally
(lower contention risk) with tri-states reserved for off-chip or bidirectional pins. Forgetting the enable decode is a classic
lab bug---two drivers fighting draws large shoot-through current.

\section{Exercises}
\label{sec:ch03-exercises}
\begin{enumerate}[label=E3.\arabic*, leftmargin=*, itemsep=0.8em]
    \item \textbf{Full adder proof.} From the truth table derive $S$ and $C_{\text{out}}$ expressions, then show $C_{\text{out}}=ab+C_{\text{in}}(a\oplus b)=ab+aC_{\text{in}}+bC_{\text{in}}$.
    \item \textbf{Ripple delay.} A FA has delay $2$ gate levels for $S$ and $1$ for $C_{\text{out}}$. What is worst-case delay of an 8-bit RCA? Where is the critical path?
    \item \textbf{MUX as logic.} Realise $F(x,y,z)=\sum m(1,2,5,6)$ using (a) an 8:1 MUX, (b) a 4:1 MUX with $x$ as data, (c) only 2:1 MUXes.
    \item \textbf{Decoder expansion.} Build a 4-to-16 decoder from 2-to-4 decoders and gates. How many 2-to-4 blocks are needed? Draw the enable wiring.
    \item \textbf{Hazard.} $F=x y + x' z$ glitches when $y=z=1$ and $x$ toggles. Show the K-map, identify the hazard, and add the consensus term $y z$ to remove it.
\end{enumerate}
% End of Chapter 3

% SC1005 --- Chapter 4: Sequential Logic --- Latches & Flip-Flops (0->100)
\chapter{Sequential Logic: Latches and Flip-Flops}
\label{ch:sequential}

\begin{quote}
\emph{``Combinational circuits forget; sequential circuits remember.''}
Add feedback and you get memory. This chapter builds memory from first principles:
SR latch $\to$ D latch $\to$ master-slave $\to$ edge-triggered flip-flop.
\end{quote}

\noindent\fbox{\parbox{\dimexpr\linewidth-2\fboxsep-2\fboxrule}{%
\textbf{From zero: memory from wires.} We assume only gates (Ch.~\ref{ch:boolean}--\ref{ch:combinational}).
We add one idea---feedback---and derive every latch/flop from it. No prior sequential knowledge.}}
\vspace{0.4em}

\section*{Learning Outcomes}
After this chapter you will be able to:
\begin{enumerate}[label=(L\arabic*)]
    \item distinguish combinational from sequential and explain the role of feedback;
    \item analyse SR and D latches built from NAND/NOR gates;
    \item explain master-slave and edge-triggered D and JK flip-flops;
    \item describe setup/hold, propagation delay, and the clock concept;
    \item use characteristic tables and excitation tables for analysis.
\end{enumerate}

% ============================================================
\section{Why Feedback? From Combinational to Sequential}
\label{sec:feedback}
A combinational output depends only on \emph{present} inputs. Add a wire from output back to
input and output now depends on \emph{history}---that is memory. The simplest feedback loop is
two cross-coupled inverters (a bistable): two stable states $0$ and $1$, holding one bit forever
until forced to flip. All latches/flops elaborate this bistable core with control inputs.

A \emph{clock} (periodic square wave) disciplines when state may change, avoiding races.
\emph{Level-sensitive} (latch) follows the input while the clock is high; \emph{edge-triggered}
(flip-flop) samples only on a clock edge ($\uparrow$ rising, $\downarrow$ falling).

\section{SR Latch: The Simplest Memory}
\label{sec:sr-latch}
Cross-coupled NOR (or NAND) gates give Set-Reset behaviour: $S=1$ sets $Q=1$, $R=1$ resets $Q=0$,
$S=R=0$ holds, $S=R=1$ is forbidden (both outputs try to be $0$, metastable when released).

\begin{figure}[h]
\centering
\begin{tikzpicture}[scale=0.80,every node/.style={font=\scriptsize}]
% SR latch NOR
\node[font=\footnotesize\bfseries] at (1.8,2.2) {SR latch (NOR)};
\node[draw,thick,fill=blue!14,minimum width=1.3cm,minimum height=0.85cm,rounded corners=2pt] (nor1) at (1.0,1.15) {NOR};
\node[draw,thick,fill=orange!16,minimum width=1.3cm,minimum height=0.85cm,rounded corners=2pt] (nor2) at (1.0,0.05) {NOR};
\draw[->,thick] (-0.45,1.35)--(0.35,1.35) node[left] {$R$};
\draw[->,thick] (-0.45,0.15)--(0.35,0.15) node[left] {$S$};
\draw[->,thick,blue!60!black] (1.65,1.15)--(2.85,1.15) node[right] {$Q$};
\draw[->,thick,red!60!black] (1.65,0.05)--(2.85,0.05) node[right] {$\bar Q$};
% feedback
\draw[->,thick,teal!70!black] (2.2,1.15)--(2.2,1.85)--(-0.15,1.85)--(-0.15,0.35)--(0.35,0.35);
\draw[->,thick,teal!70!black] (2.2,0.05)--(2.2,-0.55)--(-0.15,-0.55)--(-0.15,0.95)--(0.35,0.95);
\node[font=\tiny,teal!70!black] at (1.1,-0.85) {cross-coupled feedback};
% Truth table
\begin{scope}[xshift=4.8cm,yshift=0.15cm]
\node[font=\footnotesize\bfseries] at (1.15,1.75) {Behaviour};
\node[font=\tiny,align=left] at (1.15,1.05) {$S\;R\;|\;Q^+$};
\node[font=\tiny,align=left] at (1.15,0.65) {$0\;0\;|\;Q$ (hold)};
\node[font=\tiny,align=left] at (1.15,0.35) {$0\;1\;|\;0$ (reset)};
\node[font=\tiny,align=left] at (1.15,0.05) {$1\;0\;|\;1$ (set)};
\node[font=\tiny,align=left,fill=red!10,rounded corners=1pt] at (1.15,-0.30) {$1\;1\;|\;?$ forbidden};
\draw[thick,rounded corners=2pt] (0.05,1.35) rectangle (2.25,-0.55);
\end{scope}
\end{tikzpicture}
\caption{SR latch from cross-coupled NOR gates. Teal wires are feedback---the memory. Forbidden $S=R=1$ breaks bistability.}
\label{fig:sr-latch}
\end{figure}

Replacing NOR with NAND plus active-low inputs gives $\overline{S}\,\overline{R}$ latch,
the building block of many CMOS latches. Adding an enable $C$ gives the \emph{gated SR latch}:
when $C=0$ it holds; when $C=1$ it is transparent.

\section{D Latch: Removing the Forbidden State}
\label{sec:d-latch}
The \emph{D (data) latch} ties $R=\bar D$ so $S$ and $R$ never both $1$. One data input, one enable:
$C=1$ transparent ($Q=D$), $C=0$ hold. Solves SR's forbidden state at the cost of one lost
input combination.

\begin{figure}[h]
\centering
\begin{tikzpicture}[scale=0.80,every node/.style={font=\scriptsize}]
% D latch from SR
\node[font=\footnotesize\bfseries] at (2.0,2.15) {D latch from gated SR + inverter};
\node[draw,thick,fill=green!12,minimum width=1.0cm,minimum height=0.6cm] (inv) at (0.8,1.15) {NOT};
\draw[->,thick] (-0.3,1.15)--(0.3,1.15) node[left] {$D$};
\draw[->,thick] (1.3,1.55)--(2.1,1.55); \draw[->,thick] (1.3,0.75)--(2.1,0.75);
% SR block
\node[draw,thick,fill=blue!12,minimum width=2.0cm,minimum height=1.35cm,rounded corners=3pt] (sr) at (3.15,1.0) {Gated SR};
\draw[->,thick,red!60!black] (1.0,-0.15)--(3.15,-0.15) node[below,font=\tiny] {$C$ (enable)};
\draw[->,thick,teal!70!black] (3.15,0.15)--(3.15,-0.30)--(1.4,-0.30);
\draw[->,thick,blue!60!black] (4.15,1.15)--(5.15,1.15) node[right] {$Q$};
\draw[->,thick,red!60!black] (4.15,0.55)--(5.15,0.55) node[right] {$\bar Q$};
\node[font=\tiny,align=center] at (3.15,-0.85) {$C=1$ transparent; $C=0$ hold};
% Timing sketch
\begin{scope}[xshift=6.8cm]
\node[font=\footnotesize\bfseries] at (1.2,2.15) {Timing};
\draw[->,thick] (0,0)--(2.4,0) node[right,font=\tiny] {$t$};
\draw[thick,blue!60!black] (0,1.6)--(0.6,1.6)--(0.6,0.9)--(1.5,0.9)--(1.5,1.6)--(2.4,1.6) node[above,font=\tiny] {$D$};
\draw[thick,red!70!black] (0,0.55)--(0.9,0.55)--(0.9,1.35)--(1.8,1.35)--(1.8,0.55)--(2.4,0.55) node[below,font=\tiny] {$C$};
\draw[thick,green!60!black,dashed] (0,0.95)--(2.4,0.95) node[right,font=\tiny] {$Q$ follows $D$ when $C=1$};
\end{scope}
\end{tikzpicture}
\caption{D latch: one inverter eliminates the forbidden state. When $C=1$, $Q$ tracks $D$; when $C=0$, $Q$ holds.}
\label{fig:d-latch}
\end{figure}

\paragraph{Transparency problem.} A latch is \emph{level-sensitive}: while $C=1$, any glitch on $D$
propagates to $Q$. In multi-bit systems this causes races. The fix: make storage edge-triggered.

\section{From Latch to Flip-Flop}
\label{sec:flipflop}
\emph{Master-slave} pairs two D latches in series on opposite clock phases: master transparent
when $CLK=0$, slave when $CLK=1$---data ripples through only on the $0\!\to\!1$ edge.
Modern \emph{edge-triggered D flip-flop} ($\text{DFF}$) achieves the same with fewer transistors
(6 NAND gates) and is the standard library cell.

\begin{figure}[h]
\centering
\begin{tikzpicture}[scale=0.78,every node/.style={font=\scriptsize}]
% Master-slave
\node[font=\footnotesize\bfseries] at (2.2,2.0) {Master--slave D flip-flop};
\node[draw,thick,fill=blue!12,minimum width=1.6cm,minimum height=0.9cm,rounded corners=2pt] (m) at (1.2,0.7) {Master (D latch)};
\node[draw,thick,fill=orange!16,minimum width=1.6cm,minimum height=0.9cm,rounded corners=2pt] (s) at (3.8,0.7) {Slave (D latch)};
\draw[->,thick] (-0.2,0.85)--(0.4,0.85) node[left] {$D$};
\draw[->,thick] (2.0,0.7)--(3.0,0.7) node[midway,above,font=\tiny] {master $Q$};
\draw[->,thick,blue!60!black] (4.6,0.7)--(5.4,0.7) node[right] {$Q$};
\draw[->,thick,red!60!black] (4.6,0.3)--(5.4,0.3) node[right] {$\bar Q$};
% Clock
\draw[->,thick,teal!70!black] (1.2,-0.35)--(1.2,0.15) node[below,font=\tiny] {$\overline{CLK}$};
\draw[->,thick,teal!70!black] (3.8,-0.35)--(3.8,0.15) node[below,font=\tiny] {$CLK$};
\node[font=\tiny,fill=yellow!12,rounded corners=1pt] at (2.5,-0.85) {Samples on rising edge $\uparrow$ CLK};
% DFF symbol
\begin{scope}[xshift=7.0cm]
\node[draw,thick,fill=green!10,minimum width=1.35cm,minimum height=1.35cm] (dff) at (0.7,0.7) {};
\node[font=\tiny] at (0.15,1.0) {$D$}; \node[font=\tiny] at (0.15,0.35) {$CLK\triangleright$};
\node[font=\tiny] at (1.25,1.0) {$Q$}; \node[font=\tiny] at (1.25,0.35) {$\bar Q$};
\draw[->,thick] (-0.4,1.0)--(0.02,1.0); \draw[->,thick] (-0.4,0.35)--(0.02,0.35);
\draw[->,thick] (1.38,1.0)--(1.85,1.0); \draw[->,thick] (1.38,0.35)--(1.85,0.35);
\node[font=\footnotesize\bfseries] at (0.7,2.0) {DFF symbol};
\node[font=\tiny,fill=blue!8,rounded corners=1pt] at (0.7,-0.55) {triangle = edge-triggered};
\end{scope}
\end{tikzpicture}
\caption{Master--slave construction and the standard D flip-flop symbol (triangle denotes edge-triggered clock).}
\label{fig:master-slave}
\end{figure}

\paragraph{JK and T flip-flops.} \emph{JK} generalises SR: $J$ sets, $K$ resets, $J=K=1$ toggles
($Q^+=\bar Q$). No forbidden state. \emph{T} (toggle) has one input: $T=1$ toggles, $T=0$ holds---ideal for counters (Ch.~\ref{ch:registers}).

\paragraph{Timing constraints.} \emph{Setup time} $t_{su}$: $D$ must be stable before the edge.
\emph{Hold time} $t_h$: $D$ must remain stable after the edge. \emph{Clock-to-Q} $t_{cq}$: delay from edge to new $Q$.
Violating $t_{su}/t_h$ risks \emph{metastability} (output hovers mid-rail).

\begin{definition}[Characteristic and excitation tables]
\emph{Characteristic:} $Q^+$ as function of inputs and $Q$. D: $Q^+=D$; JK: $Q^+=JQ'+\bar KQ$.
\emph{Excitation:} inputs required to get $Q\to Q^+$. D needs $D=Q^+$; JK needs $J,K$ pattern per transition.
\end{definition}


\section{Timing Analysis and Flip-Flop Variants in Practice}
Real flip-flops have \emph{propagation delay} $t_{cq}$ (clock to $Q$), \emph{setup} $t_{su}$ and \emph{hold} $t_h$.
For a register-to-register path with combinational delay $t_{comb}$,
\[ T_{clk} \ge t_{cq}+t_{comb}+t_{su},\qquad t_{hold}\le t_{cq}+t_{comb}^{\min}. \]
If $T_{clk}$ is too short, setup fails; if $t_{comb}^{\min}$ is too short, hold fails (often fixed by buffer insertion).
\emph{Asynchronous} preset/clear override the clock---useful for power-on reset but must be deasserted synchronously to avoid metastability.

\begin{example}[Choosing a flip-flop]
D is simplest for data-path registers (one data pin, no decode). JK/T shine where toggling is the natural operation
(counters, frequency dividers: $f_{out}=f_{clk}/2^n$ from $n$ cascaded T stages). In CMOS libraries, DFF is the default cell;
JK is built as D with feedback when needed.
\end{example}

% ============================================================

% ============================================================
\section{Chapter Notes}
The bistable latch dates to Eccles--Jordan (1918) vacuum-tube trigger circuit. Edge-triggered DFFs became
standard with TTL 7474 (1970s). Setup/hold analysis is essential for static timing analysis in modern ASIC flows.

% ============================================================
\section{Metastability and Synchronizers}
When setup/hold is violated (e.g.\ an asynchronous input sampled by a clock), a flip-flop can enter \emph{metastability}---output
hovers near $V_{DD}/2$ for unbounded time before resolving to $0$ or $1$. Probability decays exponentially with extra waiting time.
Fix: \emph{two-flop synchronizer} (cascade two DFFs on the receiving clock) gives one extra cycle for resolution; mean time between
failures (MTBF) grows by orders of magnitude. Never sample an async signal with a single flop in a reliable system.

% ============================================================

\section{Exercises}
\label{sec:ch04-exercises}
\begin{enumerate}[label=E4.\arabic*, leftmargin=*, itemsep=0.8em]
    \item \textbf{SR forbidden.} For the NOR SR latch, show the $S=R=1\to0$ transition can go metastable. Why does NAND SR have the same issue with $\bar S=\bar R=0$?
    \item \textbf{D latch timing.} Sketch $C$, $D$, $Q$ for $D$ toggling while $C=1$ then holding while $C=0$. Mark where transparency lets a glitch through.
    \item \textbf{JK analysis.} Derive $Q^+=JQ'+\bar KQ$ from the JK truth table. Show $T$ is JK with $J=K=T$, and D is JK with $J=D$, $K=\bar D$.
    \item \textbf{Setup/hold.} A DFF has $t_{su}=20$\,ns, $t_h=10$\,ns, $t_{cq}=15$\,ns. Can a $50$\,MHz clock drive two DFFs in series without violation? What about $80$\,MHz?
    \item \textbf{Excitation.} Fill the excitation table for D, JK, and T flip-flops for all $Q\to Q^+$ transitions ($0\to0$, $0\to1$, $1\to0$, $1\to1$).
\end{enumerate}
% End of Chapter 4

% SC1005 --- Chapter 5: Finite State Machines (0->100)
\chapter{Finite State Machines: Mealy, Moore, and Design}
\label{ch:fsm}

\begin{quote}
\emph{``A state machine is memory with a plan.''}
Combinational logic decides; sequential logic remembers; an FSM decides \emph{based on}
what it remembers. Every controller---from vending machine to CPU---is an FSM.
\end{quote}

\noindent\fbox{\parbox{\dimexpr\linewidth-2\fboxsep-2\fboxrule}{%
\textbf{From zero: states from scratch.} We define state, transition, and output,
then derive Mealy vs.\ Moore, state diagrams, and systematic synthesis from word problems.}}
\vspace{0.4em}

\section*{Learning Outcomes}
After this chapter you will be able to:
\begin{enumerate}[label=(L\arabic*)]
    \item define FSM state, transition function, and output function;
    \item distinguish Mealy and Moore machines and convert between them;
    \item draw state diagrams and tables from word-problem specifications;
    \item synthesise synchronous FSMs using D/JK flip-flops;
    \item minimise states via partition refinement and handle unused states.
\end{enumerate}

% ============================================================
\section{What Is a Finite State Machine?}
\label{sec:fsm-def}
An FSM is a 5-tuple $(S,I,O,\delta,\lambda)$: finite state set $S$, input alphabet $I$,
output alphabet $O$, next-state function $\delta:S\times I\to S$, output function $\lambda$.
On each clock edge, state advances $s^+=\delta(s,i)$ and output is produced.

Two output styles:
\begin{itemize}
    \item \textbf{Moore:} output depends on state only, $\lambda:S\to O$ (output changes only on clock edges---glitch-free but one cycle latency).
    \item \textbf{Mealy:} output depends on state \emph{and} input, $\lambda:S\times I\to O$ (faster response, but input glitches can propagate).
\end{itemize}

\begin{figure}[h]
\centering
\begin{tikzpicture}[scale=0.82,every node/.style={font=\scriptsize}]
% Moore block
\node[draw,thick,fill=blue!12,minimum width=3.6cm,minimum height=2.0cm,rounded corners=3pt] (moore) at (2.0,0.45) {};
\node[font=\footnotesize\bfseries] at (2.0,1.75) {Moore machine};
\node[draw,thick,fill=green!14,minimum width=1.4cm,minimum height=0.7cm,rounded corners=2pt] at (1.45,0.55) {Comb.\ $\delta$};
\node[draw,thick,fill=orange!16,minimum width=1.05cm,minimum height=0.7cm] at (3.05,0.55) {State reg.};
\node[draw,thick,fill=violet!12,minimum width=1.15cm,minimum height=0.7cm,rounded corners=2pt] at (2.15,-0.55) {Output $\lambda(S)$};
\draw[->,thick] (0.15,0.55)--(0.75,0.55) node[midway,above,font=\tiny] {$I$};
\draw[->,thick,teal!70!black] (3.65,0.55)--(4.15,0.55)--(4.15,-0.55)--(2.75,-0.55);
\draw[->,thick,teal!70!black] (0.75,-0.55)--(1.55,-0.55);
\draw[->,thick] (2.15,-0.95)--(2.15,-1.35) node[below] {$O$};
\node[font=\tiny,fill=yellow!12] at (2.0,-1.65) {Output from state only};
% Mealy block
\begin{scope}[xshift=6.2cm]
\node[draw,thick,fill=blue!12,minimum width=3.6cm,minimum height=2.0cm,rounded corners=3pt] at (2.0,0.45) {};
\node[font=\footnotesize\bfseries] at (2.0,1.75) {Mealy machine};
\node[draw,thick,fill=green!14,minimum width=1.4cm,minimum height=0.7cm,rounded corners=2pt] at (1.45,0.55) {Comb.\ $\delta$};
\node[draw,thick,fill=orange!16,minimum width=1.05cm,minimum height=0.7cm] at (3.05,0.55) {State reg.};
\node[draw,thick,fill=red!12,minimum width=1.15cm,minimum height=0.7cm,rounded corners=2pt] at (2.15,-0.55) {Output $\lambda(S,I)$};
\draw[->,thick] (0.15,0.55)--(0.75,0.55) node[midway,above,font=\tiny] {$I$};
\draw[->,thick,teal!70!black] (3.65,0.55)--(4.15,0.55)--(4.15,-0.55)--(2.75,-0.55);
\draw[->,thick,teal!70!black] (0.75,-0.55)--(1.55,-0.55);
\draw[->,thick,red!60!black] (0.4,-0.15)--(1.55,-0.30) node[midway,left,font=\tiny] {$I$};
\draw[->,thick] (2.15,-0.95)--(2.15,-1.35) node[below] {$O$};
\node[font=\tiny,fill=yellow!12] at (2.0,-1.65) {Output from state + input};
\end{scope}
\end{tikzpicture}
\caption{Moore vs.\ Mealy: the only structural difference is whether input feeds the output logic (red arrow). State register is edge-triggered DFFs.}
\label{fig:moore-mealy}
\end{figure}

\section{State Diagrams and Tables}
\label{sec:state-diagram}
We specify FSMs by \emph{state diagrams} (circles = states, arcs = transitions labelled input/output)
or \emph{state tables} (rows = present state, columns = inputs). Example: sequence detector for $101$.

\begin{figure}[h]
\centering
\begin{tikzpicture}[scale=0.82,->,>=Stealth,auto,semithick,every node/.style={font=\scriptsize}]
% Moore sequence detector 101
\node[font=\footnotesize\bfseries] at (0,2.3) {Moore: detect $101$ (output $1$ in $S_3$)};
\node[state,fill=blue!12] (S0) at (0,0) {$S_0/0$};
\node[state,fill=green!14] (S1) at (3.0,0) {$S_1/0$};
\node[state,fill=orange!16] (S2) at (6.0,0) {$S_2/0$};
\node[state,fill=red!16] (S3) at (3.0,-1.85) {$S_3/1$};
\path (S0) edge[bend left=12] node {1} (S1)
      (S1) edge[bend left=12] node {0} (S2)
      (S2) edge[bend left=12] node {1} (S3)
      (S3) edge[bend left=10] node {0} (S2)
      (S0) edge[loop above] node {0} (S0)
      (S1) edge[loop above] node {1} (S1)
      (S2) edge[loop right] node {0} (S2)
      (S3) edge[loop below] node {1} (S1);
\node[font=\tiny,fill=yellow!12,rounded corners=1pt] at (3.0,-2.75) {Output is state label after `/'; glitch-free};
\end{tikzpicture}
\caption{Moore FSM that outputs $1$ when the last three inputs were $101$. State $S_i$ means ``last $i$ bits matched prefix of $101$''.}
\label{fig:seq-detector}
\end{figure}

The Mealy version needs only $3$ states (output on the arc $S_2 \xrightarrow{1/1} S_1$)
but its output can glitch if input glitches---tradeoff to choose per application.

\section{Synchronous Synthesis}
\label{sec:synthesis}
Five steps: (1) word problem $\to$ state diagram, (2) minimise states, (3) encode states in bits,
(4) derive next-state and output equations (K-maps), (5) implement with DFFs + combinational logic.

\begin{figure}[h]
\centering
\begin{tikzpicture}[scale=0.78,every node/.style={font=\scriptsize}]
% Datapath: state reg + comb logic
\node[font=\footnotesize\bfseries] at (3.0,3.2) {Synchronous FSM datapath};
\node[draw,thick,fill=orange!16,minimum width=1.4cm,minimum height=0.85cm,rounded corners=2pt,align=center] (reg) at (1.5,1.6) {State reg\\(DFFs)};
\node[draw,thick,fill=blue!12,minimum width=2.0cm,minimum height=0.85cm,rounded corners=2pt,align=center] (next) at (4.2,1.6) {Next-state\\logic};
\node[draw,thick,fill=green!12,minimum width=1.65cm,minimum height=0.7cm,rounded corners=2pt] (out) at (4.2,0.35) {Output logic};
\draw[->,thick] (-0.2,1.6)--(0.8,1.6) node[midway,above,font=\tiny] {CLK};
\draw[->,thick] (2.2,1.6)--(3.2,1.6) node[midway,above,font=\tiny] {$S$};
\draw[->,thick,red!60!black] (1.0,2.6)--(1.5,2.05) node[above,font=\tiny] {Input $I$};
\draw[->,thick,red!60!black] (1.0,2.6)--(4.2,2.05);
\draw[->,thick,teal!70!black] (4.2,1.15)--(4.2,0.70);
\draw[->,thick] (5.2,0.35)--(6.0,0.35) node[right] {$O$};
\draw[->,thick,teal!70!black] (3.2,1.15)--(1.5,0.95) node[midway,left,font=\tiny] {$S^+$};
\node[font=\tiny,fill=yellow!10,rounded corners=1pt] at (1.5,0.55) {holds $S$};
% FSM example 3 states encoded 00 01 10
\begin{scope}[xshift=7.2cm]
\node[font=\footnotesize\bfseries] at (1.3,3.2) {Example: 3-state encoding};
\node[font=\tiny] at (1.3,2.55) {$S_0=00,\;S_1=01,\;S_2=10$};
\node[font=\tiny,align=left] at (1.3,1.85) {$D_1 = S_0\cdot I$};
\node[font=\tiny,align=left] at (1.3,1.45) {$D_0 = \bar S_1 S_0 + I$};
\node[font=\tiny,align=left] at (1.3,1.05) {$O = S_1$ (Moore)};
\node[font=\tiny,fill=orange!10,rounded corners=1pt] at (1.3,0.45) {unused $11\to00$ safe};
\draw[thick,rounded corners=2pt,fill=blue!4] (-0.1,2.75) rectangle (2.7,0.15);
\end{scope}
\end{tikzpicture}
\caption{Synchronous FSM hardware: state register holds present state; combinational blocks compute next state and output. Unused codes need safe next-state.}
\label{fig:fsm-datapath}
\end{figure}

\paragraph{State minimization.} Two states are equivalent if no input sequence distinguishes them
(same outputs and transitions to equivalent states). Iteratively partition states until stable---the
minimal machine is unique up to renaming (like DFA minimization).

\paragraph{Unused states.} With $k$ bits you get $2^k$ codes; if the FSM uses fewer, the remaining
codes are don't-cares for minimization but must transition to a safe state (often reset) to avoid lock-up.


\section{State Encoding and Mealy--Moore Interconversion}
Encoding choice affects cost and speed. \emph{Binary} uses $\lceil\log_2|S|\rceil$ bits (minimal flip-flops, more decoding).
\emph{One-hot} uses $|S|$ bits with one $1$ (more flip-flops, trivial next-state logic, fast, popular in FPGAs).
\emph{Gray} minimizes bit flips between adjacent states, reducing power and glitch risk in output-coded machines.

Mealy $\to$ Moore conversion: split each state by output value (Moore state = Mealy state + output memory), at most doubling states.
Moore $\to$ Mealy: keep states, move output from state label to incoming arcs---never needs more states.
\begin{example}[Traffic-light controller]
States GREEN, YELLOW, RED with timer inputs. Moore outputs drive lamps directly (glitch-free); Mealy would combine timer expiry and state
to anticipate the next lamp---one cycle faster but needs input synchronisation to avoid lamp flicker.
\end{example}

% ============================================================

% ============================================================
\section{Chapter Notes}
FSM theory originates with Mealy (1955) and Moore (1956), both at Bell Labs. Modern synthesis tools
extract FSMs from HDL case statements and optimise encoding (binary, one-hot, Gray) automatically.

% ============================================================
\section{ASM Charts and Algorithmic State Machines}
For larger controllers, state diagrams sprawl. An \emph{ASM (Algorithmic State Machine) chart} is a flowchart with three boxes:
state box (Moore outputs), decision diamond (input test), conditional output box (Mealy outputs). Each ASM block is one state + its
transitions, making translation to HDL (one \texttt{case} per state) mechanical.

ASM example: a simple CPU fetch--decode--execute controller has states FETCH, DECODE, EXECUTE with decision diamonds for
instruction opcode and conditional outputs for register enables and ALU op codes. Timing is still synchronous---every box
executes in one clock, diamond tests are combinational within the state.

% ============================================================

\paragraph{Mealy glitch vs.\ Moore latency.} Mealy output can change mid-cycle when input glitches, so any
asynchronous input must be synchronized before feeding Mealy output logic. Moore output is glitch-free between edges
but incurs one-cycle latency (input at cycle $k$ affects output at $k+1$). Many real controllers use Mealy internally
for speed and register the output (making it Moore at the pins)---best of both worlds.

\paragraph{Timing of FSM outputs.} In a synchronous FSM, next-state is computed combinationally from present state and inputs,
then clocked into DFFs. Output logic runs in parallel. Critical path $=t_{cq}+t_{next}+t_{su}$ (Moore) or
$t_{cq}+t_{next,out}+t_{su}$ (Mealy with input-to-output). Pipelining or one-hot encoding shortens $t_{next}$ at the cost of more flops.
Static timing analysis checks this path for every state transition.

\paragraph{State reduction example.} Two states with identical outputs and transitions to equivalent next states (under current partition)
merge without changing I/O behaviour. Repeatedly refine the partition until stable---the result is the minimal machine, unique up to renaming.
Lock-up states (no input leads to the valid cycle) must be broken by adding transitions to a safe reset state.

\section{Exercises}
\label{sec:ch05-exercises}
\begin{enumerate}[label=E5.\arabic*, leftmargin=*, itemsep=0.8em]
    \item \textbf{Mealy vs.\ Moore.} Convert the Moore $101$ detector in Fig.~\ref{fig:seq-detector} to Mealy. How many states? Compare output timing on a timing diagram.
    \item \textbf{Vending machine.} Design a Moore FSM for a vending machine: inputs nickel/dime, product costs 15\textcent, output `dispense'. Draw diagram and state table.
    \item \textbf{Synthesis.} Synthesise the 3-state FSM in Fig.~\ref{fig:fsm-datapath} using D flip-flops. Write K-maps for $D_1,D_0$ and draw the gate-level circuit.
    \item \textbf{Minimization.} A FSM has 6 states with given table (provide). Apply partition refinement to find the minimal equivalent machine. Is it unique?
    \item \textbf{One-hot vs.\ binary.} Encode a 4-state FSM in binary (2 bits) and one-hot (4 bits). Compare flip-flop count, next-state logic complexity, and speed.
\end{enumerate}
% End of Chapter 5

% SC1005 --- Chapter 6: Registers and Counters (0->100)
\chapter{Registers and Counters}
\label{ch:registers}

\begin{quote}
\emph{``Registers hold; counters count; shift registers do both.''}
Built from flip-flops, these are the first useful sequential subsystems---data
storage, sequencing, and serial--parallel conversion in one chapter.
\end{quote}

\noindent\fbox{\parbox{\dimexpr\linewidth-2\fboxsep-2\fboxrule}{%
\textbf{From zero: flip-flops $\to$ systems.} We assume D/JK flip-flops (Ch.~\ref{ch:sequential})
and FSM ideas (Ch.~\ref{ch:fsm}). Each circuit is derived from ``what next-state do we want?''}}
\vspace{0.4em}

\section*{Learning Outcomes}
After this chapter you will be able to:
\begin{enumerate}[label=(L\arabic*)]
    \item build parallel registers with load and clear controls;
    \item explain shift registers and serial/parallel conversion;
    \item design ripple and synchronous counters (binary, decade, up/down);
    \item analyse counter timing, modulus, and self-correcting design;
    \item distinguish ring and Johnson counters and their decoded outputs.
\end{enumerate}

% ============================================================
\section{Registers: Parallel Storage}
\label{sec:registers}
A $n$-bit \emph{register} is $n$ D flip-flops sharing a clock and control signals.
Simplest: $n$ DFFs with common CLK---data $D_{n-1}\dots D_0$ latched on the rising edge.
Additions: \emph{asynchronous clear} ($\bar{CLR}$ forces $0$), \emph{load enable} ($LD$: when $0$, hold; when $1$, load), and tri-state outputs for bus connection.

\begin{figure}[h]
\centering
\begin{tikzpicture}[scale=0.80,every node/.style={font=\scriptsize}]
\node[font=\footnotesize\bfseries] at (2.6,2.35) {4-bit register with load and clear};
\foreach \i in {0,1,2,3} {
  \pgfmathsetmacro{\x}{\i*1.35}
  \node[draw,thick,fill=blue!\the\numexpr 12+\i*6\relax,minimum width=1.05cm,minimum height=0.85cm] at (\x+0.7,0.85) {DFF$_{\i}$};
  \draw[->,thick,blue!60!black] (\x+0.7,1.75)--(\x+0.7,1.28) node[above,font=\tiny] {$D_{\i}$};
  \draw[->,thick] (\x+0.7,0.42)--(\x+0.7,-0.1) node[below,font=\tiny] {$Q_{\i}$};
}
% shared clock and controls
\draw[->,thick,teal!70!black] (-0.5,0.55)--(0.05,0.55) node[left,font=\tiny] {CLK} -- (0.05,0.55);
\foreach \i in {0,1,2,3} { \pgfmathsetmacro{\x}{\i*1.35} \draw[thick,teal!70!black] (0.05,0.55)--(\x+0.2,0.55)--(\x+0.2,0.42); }
\draw[->,thick,red!70!black] (-0.5,1.9)--(4.9,1.9) node[right,font=\tiny] {$\overline{CLR}$};
\foreach \i in {0,1,2,3} { \pgfmathsetmacro{\x}{\i*1.35} \draw[thick,red!70!black] (\x+0.55,1.9)--(\x+0.55,1.70); }
% MUX for load
\node[font=\tiny,fill=orange!12,rounded corners=1pt] at (2.6,-0.65) {When $LD=0$, feedback holds $Q$; when $LD=1$, $D$ loads on $\uparrow$ CLK};
\node[font=\tiny,fill=yellow!10,rounded corners=1pt] at (2.6,-1.1) {$\overline{CLR}$ asynchronous: $Q\leftarrow0$ immediately};
\end{tikzpicture}
\caption{Four D flip-flops with common clock, asynchronous clear, and load-hold MUX (conceptually $D_{\text{eff}}=LD\cdot D+\overline{LD}\cdot Q$).}
\label{fig:register}
\end{figure}

Timing is identical to a single DFF: setup/hold on $D$ relative to CLK, $t_{cq}$ to $Q$.
Registers are the state elements of every FSM and data-path.

\section{Shift Registers}
\label{sec:shift}
Connect $Q_i$ to $D_{i+1}$ and data ripples one position per clock---a \emph{shift register}.
Variants: SISO, SIPO, PISO, PIPO (serial/parallel in/out). Add a MUX per stage for \emph{parallel load} vs.\ shift mode,
and a $SI$ (serial in) pin on the first stage.

Applications: serial--parallel conversion (UART), delay lines, sequence generators, and multiplication/division by $2$
(shift left/right).

\begin{figure}[h]
\centering
\begin{tikzpicture}[scale=0.80,every node/.style={font=\scriptsize}]
\node[font=\footnotesize\bfseries] at (3.2,2.0) {4-bit shift register (right shift)};
\foreach \i in {0,1,2,3} {
  \pgfmathsetmacro{\x}{\i*1.55}
  \node[draw,thick,fill=green!\the\numexpr 12+\i*7\relax,minimum width=1.20cm,minimum height=0.85cm,rounded corners=2pt] at (\x+0.8,0.7) {DFF$_{\i}$};
  \node[font=\tiny] at (\x+0.8,0.85) {bit \i};
}
\foreach \i in {0,1,2} {
  \pgfmathsetmacro{\x}{\i*1.55}
  \draw[->,thick,blue!60!black] (\x+1.40,0.7)--(\x+1.50,0.7);
}
\draw[->,thick,blue!60!black] (-0.45,0.7)--(0.20,0.7) node[left] {$SI$};
\draw[->,thick,blue!60!black] (6.40,0.7)--(6.95,0.7) node[right] {$SO$};
% clock
\draw[->,thick,teal!70!black] (3.2,-0.15)--(3.2,0.15);
\node[font=\tiny,teal!70!black] at (3.2,-0.45) {CLK (all DFFs)};
% parallel out arrows
\foreach \i in {0,1,2,3} { \pgfmathsetmacro{\x}{\i*1.55} \draw[->,thick,orange!70!black] (\x+0.8,0.25)--(\x+0.8,-0.05) node[below,font=\tiny] {$Q_{\i}$}; }
\node[font=\tiny,fill=yellow!12,rounded corners=1pt] at (3.2,-0.95) {Each $\uparrow$CLK: $Q_{i+1}\leftarrow Q_i$; $Q_0\leftarrow SI$};
% Waveform mini
\begin{scope}[xshift=8.0cm]
\node[font=\footnotesize\bfseries] at (1.1,2.0) {Shift example};
\foreach \y/\lbl in {1.35/SI,0.85/$Q_0$,0.35/$Q_1$} {
  \draw[->,thin] (0,\y)--(2.2,\y);
  \node[font=\tiny,left] at (0,\y) {\lbl};
}
\draw[thick,blue!60!black] (0,1.35)--(0.5,1.35)--(0.5,0.95)--(1.2,0.95)--(1.2,1.35)--(2.2,1.35);
\draw[thick,green!60!black] (0,0.85)--(1.0,0.85)--(1.0,0.45)--(1.7,0.45)--(1.7,0.85)--(2.2,0.85);
\node[font=\tiny] at (1.1,-0.15) {bit marches right each CLK};
\end{scope}
\end{tikzpicture}
\caption{Shift register: cascade of DFFs. Serial input enters at the left; parallel outputs tap each stage. Universal shift register adds MUX per stage for load/shift/hold.}
\label{fig:shift-reg}
\end{figure}

\section{Counters}
\label{sec:counters}
A \emph{counter} cycles through a sequence of states, usually binary counting. Modulus $m$ = number of distinct states.

\paragraph{Ripple (asynchronous) counter.} Chain T flip-flops with $Q_i$ clocking stage $i+1$.
Simple, but delay accumulates and intermediate glitches appear---unsuitable for synchronous decoding.
Mod-$2^n$ from $n$ bits; mod-$m$ ($m\neq2^n$) needs reset logic when $m$ reached.

\paragraph{Synchronous counter.} All flip-flops share CLK and toggle logic is combinational:
$T_0=1$, $T_1=Q_0$, $T_2=Q_1Q_0$, $\dots$, $T_i=\land_{j<i}Q_j$. All bits update simultaneously---faster,
no glitch, at the cost of AND gates.

\begin{figure}[h]
\centering
\begin{tikzpicture}[scale=0.80,every node/.style={font=\scriptsize}]
% Ripple 3-bit
\node[font=\footnotesize\bfseries] at (1.8,2.35) {3-bit ripple (async)};
\foreach \i in {0,1,2} {
  \pgfmathsetmacro{\x}{\i*1.45}
  \node[draw,thick,fill=blue!\the\numexpr 12+\i*8\relax,minimum width=1.05cm,minimum height=0.7cm,rounded corners=2pt] at (\x+0.7,1.15) {TFF$_{\i}$};
  \node[font=\tiny] at (\x+0.7,1.45) {$T=1$};
  \draw[->,thick] (\x+0.7,0.80)--(\x+0.7,0.35) node[below,font=\tiny] {$Q_{\i}$};
}
\draw[->,thick,teal!70!black] (-0.3,1.45)--(0.15,1.45) node[left,font=\tiny] {CLK};
\draw[->,thick,red!60!black] (1.23,0.55)--(1.60,1.05); \draw[->,thick,red!60!black] (2.68,0.55)--(3.05,1.05);
\node[font=\tiny,fill=red!10] at (1.8,0.0) {$Q_0$ clocks $Q_1$, etc.};
% Synchronous 3-bit
\begin{scope}[xshift=5.5cm]
\node[font=\footnotesize\bfseries] at (1.8,2.35) {3-bit synchronous};
\foreach \i in {0,1,2} {
  \pgfmathsetmacro{\x}{\i*1.25}
  \node[draw,thick,fill=orange!\the\numexpr 12+\i*8\relax,minimum width=0.95cm,minimum height=0.7cm,rounded corners=2pt] at (\x+0.7,1.15) {TFF$_{\i}$};
  \draw[->,thick] (\x+0.7,0.80)--(\x+0.7,0.35) node[below,font=\tiny] {$Q_{\i}$};
}
\draw[->,thick,teal!70!black] (0.2,0.10)--(0.70,0.55); \draw[->,thick] (1.25,0.10)--(1.95,0.55); \draw[->,thick] (2.50,0.10)--(3.20,0.55);
\node[font=\tiny,fill=green!10,rounded corners=1pt] at (1.8,1.80) {$T_0=1,\;T_1=Q_0,\;T_2=Q_1Q_0$};
\draw[->,thick,teal!70!black] (-0.4,0.10)--(1.8,0.10) node[below,font=\tiny] {CLK (common)};
\node[font=\tiny,fill=yellow!12] at (1.8,0.0) {All update together};
\end{scope}
\end{tikzpicture}
\caption{Ripple vs.\ synchronous 3-bit binary counter. Ripple uses the previous $Q$ as clock (red); synchronous shares CLK and gates the $T$ inputs.}
\label{fig:counters}
\end{figure}

\paragraph{Up/down and decade.} Add a direction bit to select increment vs.\ decrement.
Decade ($0$--$9$) and other non-power-of-two moduli decode terminal count and synchronously reset or use
natural BCD behaviour (e.g.\ 74LS90).

\paragraph{Ring and Johnson.} \emph{Ring}: shift register with $Q_{n-1}$ fed back to $D_0$---one hot circulates ($n$ states, simple decode).
\emph{Johnson (twisted ring)}: feed $\bar Q_{n-1}$ back---$2n$ states with Gray-like adjacency, decoded by 2-literal ANDs.


\section{Counter Modulus, Self-Correction, and Timing}
Modulus $m$ need not be a power of two. For synchronous mod-$m$ ($m<2^n$), decode terminal count $m-1$ and synchronously
load $0$ (or use natural truncated sequence with don't-cares for $m$--$2^n-1$ to minimize logic). Asynchronous reset for mod-$m$
creates a short glitch state (e.g.\ $101\to000$ via $101$ spike)---use synchronous reset in clocked systems.

\emph{Self-correction:} every unused state must eventually enter the valid cycle. Verify by tracing next-state from each unused code;
if a lock-up cycle exists, add logic to break it (usually force to $0$ on next clock). Ring/Johnson counters are inherently
self-correcting only if designed with correction gates (e.g.\ start-up reset).
\begin{example}[Mod-10 decade]
BCD decade $0$--$9$ from 4 bits: states $1010$--$1111$ are don't-cares. K-maps for next-state treat them as $\times$,
yielding simpler logic than a full 16-state map. Add synchronous clear when $1001\to0000$ to close the cycle.
\end{example}

% ============================================================

% ============================================================
\section{Chapter Notes}
Counters and shift registers were the first MSI building blocks (7495, 74193, 74164).
Pipelined data-paths still use registers as pipeline boundaries; Johnson counters appear in stepper-motor drivers.

% ============================================================
\section{Applications: From Counters to Timers}
Counters plus registers build timers, frequency dividers, and UART baud generators. A down-counter loaded with $N$ and counting to $0$
produces a pulse every $N+1$ clocks (programmable period). A prescaler (power-of-two divider) followed by a programmable counter
covers wide timing ranges with modest bit-width. In displays, a BCD counter plus 7-segment decoder drives each digit, multiplexed
at $\sim1$\,kHz so persistence of vision shows a stable number.

% ============================================================

\paragraph{Shift-register applications.} Beyond storage, shift registers implement \emph{sequence generators}
(LFSR with XOR feedback gives maximal $2^n-1$ pseudo-random sequence), \emph{CRC} computation (polynomial division
via shifted XOR), and \emph{serial protocols} (SPI/I$^2$C shift data on SCLK). The universal shift register's four
modes (hold, shift-left, shift-right, parallel load) are selected by two mode bits---one chip, many uses.

\section{Exercises}
\label{sec:ch06-exercises}
\begin{enumerate}[label=E6.\arabic*, leftmargin=*, itemsep=0.8em]
    \item \textbf{Register timing.} A 4-bit register uses DFFs with $t_{su}=15$\,ns, $t_{cq}=20$\,ns. What is max CLK frequency? What changes with a load MUX adding 8\,ns before $D$?
    \item \textbf{Shift modes.} Draw a 4-bit universal shift register stage with MUX selecting hold/shift-left/shift-right/parallel-load. Give the $S_1S_0$ table.
    \item \textbf{Synchronous design.} Design a mod-6 synchronous binary counter (0--5) using T flip-flops. Derive $T_2T_1T_0$ via K-maps and handle states 6,7.
    \item \textbf{Ripple glitch.} Show a 3-bit ripple counter glitch when $011\to100$: list transient binary values and explain why synchronous avoids them.
    \item \textbf{Johnson decode.} For a 4-bit Johnson counter, list its $8$ states and give minimal decode for each state using 2-input gates.
\end{enumerate}
% End of Chapter 6

% SC1005 --- Chapter 7: Arithmetic Circuits (0->100)
\chapter{Arithmetic Circuits: Fast Adders, Multipliers, and ALUs}
\label{ch:arithmetic}

\begin{quote}
\emph{``Addition is the mother of all arithmetic.''}
Fast addition, subtraction, multiplication, and the ALU that unites them---all
derived from the full adder and the idea of \emph{anticipating} the carry.
\end{quote}

\noindent\fbox{\parbox{\dimexpr\linewidth-2\fboxsep-2\fboxrule}{%
\textbf{From zero: arithmetic from gates.} We assume adders and two's complement
(Ch.~\ref{ch:numbers}--\ref{ch:combinational}). Each faster circuit is motivated by the delay
of the previous one, then derived and costed.}}
\vspace{0.4em}

\section*{Learning Outcomes}
After this chapter you will be able to:
\begin{enumerate}[label=(L\arabic*)]
    \item compute generate/propagate and derive carry-lookahead equations;
    \item design carry-lookahead and carry-select adders and compare delay vs.\ cost;
    \item build array and shift-add multipliers and explain Booth recoding;
    \item perform two's complement subtraction and overflow detection in hardware;
    \item describe the ALU organization and status flags ($C,Z,N,V$).
\end{enumerate}

% ============================================================
\section{Why Ripple Is Too Slow}
\label{sec:ripple-slow}
A $32$-bit ripple-carry adder waits for carry to walk $32$ stages---$\sim64$ gate delays.
At $3$\,GHz that is several clock cycles. Faster adders \emph{compute} carries in parallel
instead of waiting for them.

Define for each bit $i$: generate $g_i=a_i b_i$ (creates a carry), propagate $p_i=a_i\oplus b_i$
(passes an incoming carry), and $c_{i+1}=g_i+p_i c_i$. This recurrence is the key.

\section{Carry-Lookahead Adder (CLA)}
\label{sec:cla}
Unroll $c_{i+1}=g_i+p_i c_i$ to get sum-of-products for each carry in terms of $g,p$ and $c_0$:
\begin{align*}
c_1 &= g_0+p_0c_0,\\
c_2 &= g_1+p_1g_0+p_1p_0c_0,\\
c_3 &= g_2+p_2g_1+p_2p_1g_0+p_2p_1p_0c_0,\quad\ldots
\end{align*}
Each $c_i$ is two gate levels (AND-OR) after $g,p$ are ready---$O(1)$ depth, independent of $n$,
at the cost of fan-in. Practical CLAs are grouped (4-bit blocks with block $G,P$) and cascaded.

\begin{figure}[h]
\centering
\begin{tikzpicture}[scale=0.78,every node/.style={font=\scriptsize}]
\node[font=\footnotesize\bfseries] at (3.8,3.15) {4-bit carry-lookahead: $g_i,p_i$ $\to$ lookahead $\to$ sums};
% g/p gen
\node[draw,thick,fill=blue!10,minimum width=1.35cm,minimum height=1.35cm,rounded corners=2pt] (gp) at (1.35,1.55) {$g,p$ gen};
\node[font=\tiny] at (1.35,1.0) {$g_i=a_i b_i$};
\node[font=\tiny] at (1.35,0.70) {$p_i=a_i\oplus b_i$};
\foreach \i in {0,1,2,3} \draw[->,thin,blue!60!black] (0.15,2.15-\i*0.28)--(0.65,2.15-\i*0.28) node[left,font=\tiny] {$a_{\i}b_{\i}$};
% Lookahead block
\node[draw,thick,fill=green!12,minimum width=2.15cm,minimum height=1.65cm,rounded corners=2pt] (la) at (3.95,1.55) {Lookahead};
\node[font=\tiny] at (3.95,1.05) {$c_{i+1}=g_i+p_i c_i$};
\node[font=\tiny,fill=yellow!12] at (3.95,0.70) {AND-OR, 2 levels};
\draw[->,thick,orange!70!black] (2.02,1.55)--(2.87,1.55) node[midway,above,font=\tiny] {$g,p$};
\draw[->,thick,red!60!black] (1.35,2.65)--(3.95,2.45) node[midway,above,font=\tiny] {$c_0$};
\foreach \i in {1,2,3,4} \draw[->,thick,teal!70!black] (3.95,2.35-\i*0.28)--(6.3,2.35-\i*0.28) node[right,font=\tiny] {$c_{\i}$};
% Sums
\node[draw,thick,fill=orange!14,minimum width=1.45cm,minimum height=1.65cm,rounded corners=2pt] (sum) at (6.65,1.55) {XOR sums};
\node[font=\tiny] at (6.65,0.95) {$S_i=p_i\oplus c_i$};
\draw[->,thick,blue!60!black] (5.02,1.10)--(5.92,1.10) node[midway,above,font=\tiny] {$p$};
\draw[->,thick] (6.65,0.35)--(6.65,-0.25) node[below,font=\tiny] {$S_3\dots S_0$};
\node[font=\tiny,fill=violet!10,rounded corners=1pt] at (3.95,-0.55) {Critical path: $g,p$ (2) + lookahead (2) + XOR (1) $\approx$ 5 levels};
\end{tikzpicture}
\caption{Carry-lookahead dataflow. All carries are computed in parallel from $g,p,c_0$ in two AND-OR levels, then XORed for sums.}
\label{fig:cla}
\end{figure}

\emph{Carry-select} trades off: compute two speculative results ($c_{\text{mid}}=0$ and $1$) for the upper half
and MUX-select when the real carry arrives---less logic than full CLA, faster than ripple.

\section{Subtraction and Overflow in Hardware}
Subtraction $A-B$ is $A+\bar B+1$ (two's complement negate: invert $B$ and set $c_0=1$). One adder does both
add and subtract with a $\mathit{SUB}$ control that XORs $B$ with $\mathit{SUB}$ and feeds $\mathit{SUB}$ as $c_0$.
Overflow $V = c_n\oplus c_{n-1}$ (carry into sign $\neq$ carry out) or equivalently, same-sign inputs producing opposite-sign sum.

\section{Multiplication}
\label{sec:multiplication}
Paper-and-pencil binary multiplication is shifts and adds. Two hardware styles:

\begin{itemize}
    \item \textbf{Array multiplier:} $n\times n$ grid of AND gates (partial products) plus adders summing rows---$O(n^2)$ gates, delay $O(n)$.
    \item \textbf{Shift-add (sequential):} one adder + shift register, iterating $n$ cycles---compact, used in microcode.
\end{itemize}
\emph{Booth recoding} handles signed two's complement directly by recoding runs of $1$s (e.g.\ $01110 = +16-2$),
reducing partial products and speeding signed multiply.

\begin{figure}[h]
\centering
\begin{tikzpicture}[scale=0.78,every node/.style={font=\scriptsize}]
\node[font=\footnotesize\bfseries] at (2.8,3.0) {4$\times$4 array multiplier (partial products + adder rows)};
% Partial products grid
\foreach \i in {0,1,2,3} {
  \foreach \j in {0,1,2,3} {
    \pgfmathsetmacro{\x}{1.0+\j*0.62}
    \pgfmathsetmacro{\y}{2.15-\i*0.48}
    \node[draw,fill=blue!\the\numexpr 10+\i*6\relax,minimum width=0.50cm,minimum height=0.32cm,font=\tiny] at (\x,\y) {$a_{\j}b_{\i}$};
  }
  \node[font=\tiny,red!60!black] at (3.95,2.15-\i*0.48) {$\leftarrow b_{\i}$};
}
\foreach \j in {0,1,2,3} \node[font=\tiny,blue!60!black] at (1.0+\j*0.62,2.55) {$a_{\j}$};
% Shift illustration
\draw[->,thick,orange!70!black] (1.0,1.60)--(0.15,1.15) node[left,font=\tiny] {row $i$ shifted $i$};
% Adder rows
\node[draw,thick,fill=green!12,minimum width=3.2cm,minimum height=0.45cm,rounded corners=2pt] at (2.2,0.95) {Adder row 1};
\node[draw,thick,fill=orange!14,minimum width=3.2cm,minimum height=0.45cm,rounded corners=2pt] at (2.2,0.45) {Adder row 2};
\node[draw,thick,fill=violet!12,minimum width=3.2cm,minimum height=0.45cm,rounded corners=2pt] at (2.2,-0.05) {Adder row 3};
\draw[->,thick] (2.2,-0.32)--(2.2,-0.65) node[below,font=\tiny] {$P_7\dots P_0$ (8-bit product)};
\node[font=\tiny,fill=yellow!12] at (2.2,-1.05) {Each row is a ripple/CLA adder summing one more partial product};
\end{tikzpicture}
\caption{Array multiplier: each row of AND gates creates one partial product; successive adder rows accumulate them with increasing shift.}
\label{fig:array-mult}
\end{figure}

\section{Arithmetic Logic Unit (ALU)}
The ALU is a MUX of operations sharing one adder: inputs $A,B$, control $Op$ selects
arithmetic ($+,-,$ increment) or logic (AND, OR, XOR, NOT) plus shifts. Status flags:
\emph{C} carry out, \emph{Z} zero ($NOR$ of result), \emph{N} sign (MSB), \emph{V} signed overflow, used by branch logic.

\begin{figure}[h]
\centering
\begin{tikzpicture}[scale=0.78,every node/.style={font=\scriptsize}]
\node[draw,thick,fill=blue!10,minimum width=3.4cm,minimum height=2.15cm,rounded corners=3pt] (alu) at (2.5,0.75) {ALU};
\node[font=\footnotesize\bfseries] at (2.5,2.15) {ALU ($Op$ selects function)};
\draw[->,thick,blue!60!black] (0.15,1.55)--(0.80,1.55) node[left] {$A$};
\draw[->,thick,blue!60!black] (0.15,0.75)--(0.80,0.75) node[left] {$B$};
\draw[->,thick,red!70!black] (2.5,2.45)--(2.5,1.82) node[above,font=\tiny] {$Op$ (3--4 bits)};
\draw[->,thick] (4.20,1.15)--(5.35,1.15) node[right] {$F$ (result)};
% flags
\node[draw,thick,fill=orange!14,minimum width=1.05cm,minimum height=0.45cm,rounded corners=1pt] at (5.05,0.25) {Flags};
\node[font=\tiny] at (5.05,0.25) {$C\,Z\,N\,V$};
\draw[->,thick,teal!70!black] (4.00,0.25)--(4.50,0.25);
\node[font=\tiny,align=left] at (2.5,1.05) {ADD, SUB, AND, OR,\\XOR, NOT, SLL, SRL};
\node[font=\tiny,fill=yellow!12] at (2.5,-0.55) {Adder + logic + shifter + output MUX};
\end{tikzpicture}
\caption{ALU organization: one shared adder plus logic and shifter, multiplexed by the operation code; flags summarize the result for branching.}
\label{fig:alu}
\end{figure}


\section{Carry-Select, Prefix Adders, and Multiplication Tradeoffs}
\emph{Carry-select} partitions $n$ bits into blocks (e.g.\ $4+4+4+4$ for $n=16$), precomputes both $c_{in}=0$ and $1$ per block,
then MUX-selects. Delay $\approx$ block adder + $(n/{\rm blocksize}-1)$ MUX delays---between ripple and full CLA.
\emph{Parallel-prefix} (Kogge--Stone, Brent--Kung) generalises CLA to $O(\log n)$ depth with $O(n\log n)$ or $O(n)$ cost,
the standard choice for 32/64-bit ALU adders at GHz frequencies.

Multiplication area--time tradeoff: array is $O(n^2)$ gates but single-cycle; shift-add is $O(n)$ gates but $n$ cycles.
For $n=32$, Wallace/Dadda trees compress partial products with 3:2 counter arrays to $O(\log n)$ stages, used in high-performance multipliers.
\begin{example}[Flag generation]
Result $R=A+B$: $C=$ carry out, $Z=\bigwedge_i \bar R_i$ (NOR), $N=R_{n-1}$, $V=C_n\oplus C_{n-1}$.
Branch conditions combine them: e.g.\ signed $A<B$ iff $N\oplus V=1$ for two's complement.
\end{example}

% ============================================================

% ============================================================
\section{Chapter Notes}
CLA was described by Weinberger and Smith (1958); Booth's algorithm (1951) sped early computer multipliers.
Modern ALUs use parallel-prefix adders (Kogge--Stone, Brent--Kung) that generalise CLA to $O(\log n)$ delay.

% ============================================================
\section{Division Preview and Fixed-Point Scaling}
Division is the hardest basic operation---iterative subtract-and-shift (restoring or non-restoring) takes $n$ cycles for $n$ bits,
with careful handling of two's complement signs. We omit full design here (computer arithmetic course), but note that multiplication
and division by $2^k$ are just shifts, and fixed-point scaling (e.g.\ Q15 format: 1 sign + 15 fractional bits) uses this to implement
fractional arithmetic with integer hardware, ubiquitous in DSP and embedded control.

% ============================================================

\paragraph{Signed overflow vs.\ carry.} Carry $C$ is unsigned overflow; $V$ is signed overflow. For addition,
$V = (A_{n-1}\!\oplus\!R_{n-1})\cdot(B_{n-1}\!\oplus\!R_{n-1})$ or $C_n\!\oplus\!C_{n-1}$. Hardware computes both
in parallel from the adder's internal carries. Branch logic selects which flag to test based on signed vs.\ unsigned
interpretation of the preceding operation (e.g.\ \texttt{BCC} vs.\ \texttt{BVS} in 6502/ARM).

\paragraph{Shifter and rotator.} A \emph{barrel shifter} rotates or shifts $n$ bits by any amount in $O(\log n)$ MUX stages
(e.g.\ 32-bit shift in 5 stages of 2:1 MUXes controlled by shift-amount bits). Shifts are essential for normalization in
floating-point and for alignment in addition. Rotates preserve bits (useful in cryptography); arithmetic right shift replicates
the sign bit to preserve two's complement value.

\paragraph{Multiplication correctness sketch.} Array multiplier sums partial products $a_jb_i$ at column $j+i$; each adder row preserves
$\sum$ partials = product. By induction on rows, after $k$ rows the partial sum equals $\sum_{i<k}\sum_j a_jb_i2^{j+i}$,
and after $n$ rows the full $2n$-bit product. Booth's recoding maintains the same invariant with fewer non-zero partials.

\paragraph{Cost vs.\ delay summary.} Ripple: $O(n)$ gates, $O(n)$ delay. CLA: $O(n^2)$ (na\"ive) or $O(n)$ with grouping, $O(\log n)$ delay. Carry-select: $\approx1.5\times$ ripple gates, $\approx O(\sqrt n)$ delay. Choose by constraints: area-limited $\to$ ripple/carry-select, speed-limited $\to$ prefix CLA.

\section{Exercises}
\label{sec:ch07-exercises}
\begin{enumerate}[label=E7.\arabic*, leftmargin=*, itemsep=0.8em]
    \item \textbf{CLA equations.} For $n=4$, write $c_1\dots c_4$ in terms of $g_i,p_i,c_0$. Group as 4-bit block with $G=g_3+p_3g_2+p_3p_2g_1+p_3p_2p_1g_0$, $P=p_3p_2p_1p_0$.
    \item \textbf{Delay.} Ripple $n=16$ with $t_{FA}=2$ vs.\ CLA with $t_{gp}=2$, $t_{look}=2$, $t_{xor}=1$. Compute worst-case delays and speedup.
    \item \textbf{Subtractor.} Show how one adder with $SUB$ control computes both $A+B$ and $A-B$. Where does overflow $V$ come from? Test $A=0111$, $B=0001$.
    \item \textbf{Booth.} Multiply $6\times(-3)$ ($0101\times1101$ in 4-bit two's complement) using Booth recoding. List partial products and final product.
    \item \textbf{ALU flags.} For $A=1011$, $B=0101$, compute $A+B$, $A-B$, $A\;\&\;B$ in 4-bit two's complement and give $C,Z,N,V$ for each.
\end{enumerate}
% End of Chapter 7

% SC1005 --- Chapter 8: Memory and Programmable Logic (0->100)
\chapter{Memory and Programmable Logic}
\label{ch:memory}

\begin{quote}
\emph{``Logic computes; memory remembers across power cycles and clock cycles.''}
From one-bit cells to RAM, ROM, and the programmable fabrics that let you
\emph{draw} a circuit into existence---the storage side of digital design.
\end{quote}

\noindent\fbox{\parbox{\dimexpr\linewidth-2\fboxsep-2\fboxrule}{%
\textbf{From zero: memory from latches.} We build up from latch/register (Ch.~\ref{ch:registers})
to arrays, then show how any truth table becomes a memory lookup, and how FPGAs generalise this.}}
\vspace{0.4em}

\section*{Learning Outcomes}
After this chapter you will be able to:
\begin{enumerate}[label=(L\arabic*)]
    \item explain SRAM and DRAM cells and their tradeoffs;
    \item describe address decoding, read/write timing, and memory expansion;
    \item distinguish ROM, PROM, EPROM, EEPROM, and Flash;
    \item use PLA, PAL, and ROM as two-level logic implementations;
    \item describe FPGA organization (LUTs, interconnect, configuration).
\end{enumerate}

% ============================================================
\section{Memory Organization}
\label{sec:memory-org}
A $2^n\times m$ memory holds $2^n$ words of $m$ bits. Inputs: $n$ address lines select one word,
$m$ data lines carry the word, plus $\overline{CS}$ (chip select), $R/\bar W$ (read vs.\ write),
and $\overline{OE}$ (output enable). Internally: a matrix of cells, row/column decoders, and sense amplifiers.

Capacity examples: $1\text{K}\times8 = 8192$ bits $=1$\,KB. Address space grows exponentially with $n$;
widening $m$ replicates arrays in parallel.

\begin{figure}[h]
\centering
\begin{tikzpicture}[scale=0.80,every node/.style={font=\scriptsize}]
\node[font=\footnotesize\bfseries] at (3.0,3.0) {$2^n\times m$ memory organization};
% Memory array
\node[draw,thick,fill=blue!8,minimum width=3.2cm,minimum height=1.85cm] (array) at (2.4,1.15) {Cell array};
\foreach \y in {0.45,0.85,1.25,1.65} \draw[thin,gray] (0.8,\y)--(4.0,\y);
\foreach \x in {1.15,1.75,2.35,2.95,3.55} \draw[thin,gray] (\x,0.22)--(\x,2.08);
\node[font=\tiny,fill=green!10,rounded corners=1pt] at (2.4,1.70) {$2^n$ rows $\times$ $m$ cols};
% Row decoder
\node[draw,thick,fill=orange!14,minimum width=1.05cm,minimum height=1.85cm,rounded corners=2pt,align=center] (rdec) at (0.05,1.15) {Row\\dec.};
\draw[->,thick,blue!60!black] (-1.05,1.55)--(-0.48,1.55) node[left,font=\tiny] {Addr};
\draw[->,thick,teal!70!black] (0.57,1.15)--(0.80,1.15);
\node[font=\tiny] at (0.55,0.45) {$n$ bits};
% Data I/O
\node[draw,thick,fill=violet!12,minimum width=3.2cm,minimum height=0.45cm,rounded corners=2pt] (io) at (2.4,0.05) {I/O / sense amps};
\draw[->,thick] (2.4,-0.55)--(2.4,-0.18) node[below,font=\tiny] {Data $m$ bits};
\draw[<->,thick,red!60!black] (4.35,1.15)--(4.90,1.15) node[right,font=\tiny] {$R/\bar W$};
% Expansion hint
\node[font=\tiny,fill=yellow!12,rounded corners=1pt] at (2.4,-1.05) {Row dec. selects one word; column MUX selects bits; $R/\bar W$ steers direction};
\end{tikzpicture}
\caption{Memory as a matrix: address decoded to one row, data flows through sense amplifiers. Controls $R/\bar W$, $\overline{CS}$, $\overline{OE}$ gate the transfer.}
\label{fig:memory-org}
\end{figure}

\section{SRAM vs.\ DRAM Cell}
\label{sec:sram-dram}
\emph{SRAM} cell: 6-transistor bistable (cross-coupled inverters + access transistors)---fast, no refresh, larger area.
\emph{DRAM} cell: 1 transistor + 1 capacitor---tiny, dense, but charge leaks and must be refreshed every few ms,
and reads are destructive (charge sharing) so a write-back is required. SRAM for caches, DRAM for main memory.

ROM variants: mask ROM (factory), PROM (one-time fuse), EPROM (UV-erase), EEPROM/Flash (electrically erase in blocks)---non-volatile.

\section{Programmable Logic: PLA, PAL, ROM}
Any $n$-input, $m$-output combinational function is a $2^n\times m$ lookup table: address = input code,
data = output values. Using a ROM this way trades silicon for universality but wastes capacity when don't-cares abound.

\emph{PLA} (Programmable Logic Array): programmable AND plane (product terms) followed by programmable OR plane---
minimal two-level realization. \emph{PAL} (Programmable Array Logic): programmable AND, fixed OR---cheaper,
faster, slightly less flexible. Both are field-programmable (fuses/antifuses/Flash).

\begin{figure}[h]
\centering
\begin{tikzpicture}[scale=0.80,every node/.style={font=\scriptsize}]
\node[font=\footnotesize\bfseries] at (3.0,2.85) {PLA vs.\ ROM as logic (3 inputs, 2 outputs)};
% PLA
\begin{scope}
\node[draw,thick,fill=blue!10,minimum width=2.15cm,minimum height=1.05cm,rounded corners=2pt] (and) at (1.35,1.55) {AND plane};
\node[font=\tiny] at (1.35,1.25) {product terms};
\node[draw,thick,fill=green!12,minimum width=2.15cm,minimum height=0.85cm,rounded corners=2pt] (or) at (1.35,0.55) {OR plane};
\draw[->,thick,blue!60!black] (0.15,1.55)--(0.27,1.55) node[left,font=\tiny] {$x,y,z$};
\draw[->,thick,teal!70!black] (1.35,1.02)--(1.35,0.97);
\draw[->,thick] (1.35,0.12)--(1.35,-0.35) node[below,font=\tiny] {$F_1,F_2$};
\node[font=\tiny,fill=yellow!10] at (1.35,-0.75) {PLA: only needed terms};
\end{scope}
% ROM
\begin{scope}[xshift=4.2cm]
\node[draw,thick,fill=orange!14,minimum width=2.15cm,minimum height=1.35cm,rounded corners=2pt] (rom) at (1.35,1.35) {ROM $2^n\times m$};
\node[font=\tiny] at (1.35,1.05) {$2^n$ words};
\draw[->,thick,blue!60!black] (0.15,1.55)--(0.27,1.55) node[left,font=\tiny] {$x,y,z$};
\node[font=\tiny] at (0.05,1.85) {addr};
\draw[->,thick] (1.35,0.67)--(1.35,-0.35) node[below,font=\tiny] {$F_1,F_2$};
\node[font=\tiny] at (1.55,-0.75) {data};
\node[font=\tiny,fill=red!10] at (1.35,-1.15) {ROM: all $2^n$ combos};
\end{scope}
% Efficiency note
\node[font=\tiny,align=center,fill=violet!8,rounded corners=2pt] at (3.0,-1.65) {PLA $\ll$ ROM when few product terms; ROM wins for arbitrary/many functions};
\end{tikzpicture}
\caption{Two ways to make any combinational function: PLA creates only the product terms you need; ROM stores the full truth table.}
\label{fig:pla-rom}
\end{figure}

\section{FPGAs: Logic You Configure}
\label{sec:fpga}
A Field-Programmable Gate Array tiles \emph{Configurable Logic Blocks} (CLBs) containing
Look-Up Tables (LUTs, typically 4--6 input SRAM memories that \emph{are} tiny ROMs), flip-flops,
and MUXes, stitched by a programmable interconnect fabric. Configuration is a bitstream loaded
from Flash/SRAM that sets every LUT truth table and every routing switch.

Design flow: HDL $\to$ synthesis $\to$ mapping to LUTs $\to$ place \& route $\to$ bitstream.
Reprogrammable in milliseconds---prototyping, acceleration, and low-volume production.

\begin{figure}[h]
\centering
\begin{tikzpicture}[scale=0.78,every node/.style={font=\scriptsize}]
\node[font=\footnotesize\bfseries] at (3.5,2.85) {FPGA fabric (CLBs + interconnect)};
\foreach \x in {0,1,2,3} {
  \foreach \y in {0,1,2} {
    \pgfmathsetmacro{\px}{\x*1.55}
    \pgfmathsetmacro{\py}{\y*0.62+0.15}
    \node[draw,thick,fill=blue!\the\numexpr 10+\y*5\relax,minimum width=1.15cm,minimum height=0.48cm,rounded corners=1pt] at (\px+0.68,\py) {CLB};
    \node[font=\tiny] at (\px+0.68,\py) {LUT+FF};
  }
}
% Interconnect channels
\draw[thick,teal!60!black,dashed] (-0.15,0.0)--(6.55,0.0) node[right,font=\tiny] {routing channel};
\draw[thick,teal!60!black,dashed] (-0.15,0.62)--(6.55,0.62);
\draw[thick,teal!60!black,dashed] (-0.15,1.24)--(6.55,1.24);
% Labels
\node[font=\tiny,fill=green!10,rounded corners=1pt] at (1.2,-0.55) {LUT = $k$-input SRAM truth table};
\node[font=\tiny,fill=orange!10,rounded corners=1pt] at (4.5,-0.55) {Interconnect = MUX switches};
\node[font=\tiny,fill=yellow!12] at (3.5,-0.95) {Bitstream sets every LUT and switch};
% I/O ring
\draw[thick,rounded corners=3pt,fill=violet!6] (-0.35,-0.20) rectangle (6.85,2.15);
\node[font=\tiny,violet!70!black] at (3.5,2.45) {I/O blocks around perimeter};
\end{tikzpicture}
\caption{FPGA: array of CLBs (each a LUT plus optional flip-flop) connected by a programmable routing fabric. The bitstream configures LUT contents and interconnect.}
\label{fig:fpga}
\end{figure}

Memory hierarchy note: registers (fastest, fewest), SRAM caches, DRAM main memory, Flash/SSD storage---each level trades speed for capacity and cost. Temporal and spatial locality make the hierarchy effective.


\section{Read/Write Timing and Memory Expansion}
Read cycle: place address, assert $\overline{CS}$ and $\overline{OE}$, wait $t_{AA}$ (address access time), sample data.
Write cycle: place address and data, pulse $\overline{WE}$ or $R/\bar W$, meet setup/hold $t_{AS},t_{DS},t_{DH}$. SRAM is asynchronous
or synchronous (clocked, pipelined for cache). DRAM needs $\overline{RAS}/\overline{CAS}$ multiplexing and periodic refresh
(burst or distributed, $\sim64$\,ms retention).

Expansion: widening ($m$) parallels chips sharing address; deepening ($2^n$) decodes upper address bits to chip-selects.
Example: $4\text{K}\times8$ from $1\text{K}\times8$ needs $4$ chips, $2$ address bits into a 2-to-4 decoder whose outputs are $\overline{CS}_i$,
data buses tied together (tri-state or wired-OR with $\overline{OE}$ gating).
\begin{example}[ROM as lookup]
A $32\times8$ ROM stores ASCII codes: address $=5$-bit input code, data $=8$-bit ASCII. Changing the function is a ROM reprogram,
not a rewiring---the original ``firmware'' idea before Flash and FPGAs made reconfiguration routine.
\end{example}

% ============================================================

% ============================================================
\section{Chapter Notes}
Wilkes (EDSAC, 1949) introduced stored-program memory; Dennard's 1T DRAM (1968) enabled cheap main memory.
PLAs/PALs dominated glue logic in the 1980s; Xilinx XC2064 (1985) launched the FPGA era (now AMD/Xilinx and Intel/Altera).

% ============================================================
\section{Memory Hierarchy and Locality}
Registers ($t\sim0.3$\,ns) $\to$ SRAM cache ($1$--$3$\,ns) $\to$ DRAM ($50$--$100$\,ns) $\to$ Flash/SSD ($50\,\mu$s--$1$\,ms) $\to$ disk.
Each level is larger, cheaper per bit, and slower. \emph{Temporal locality} (reuse) and \emph{spatial locality} (nearby addresses)
make caches effective: a small fast cache captures most accesses. Write policies (write-through vs.\ write-back) and replacement
(LRU) manage coherence between levels---full treatment in computer architecture (SC1007/SC2002).

% ============================================================

\paragraph{Error detection.} Memories add parity or ECC: single-bit parity detects one error; Hamming $(72,64)$
corrects one bit and detects two (SECDED) with 8 check bits per 64-bit word---standard in servers. Flash adds wear
leveling and bad-block management since cells degrade after $\sim10^5$ program/erase cycles; the controller remaps
writes to spread wear, turning the memory into a small computer itself.

\paragraph{Non-volatile evolution.} Mask ROM gave way to PROM (fusible links), EPROM (floating-gate, UV erase), EEPROM
(byte-erase via Fowler--Nordheim tunneling), and Flash (block-erase, NOR for XIP vs.\ NAND for density). Emerging NVM---MRAM,
ReRAM, PCM---promise DRAM-like speed with non-volatility, potentially collapsing the memory hierarchy in future architectures.
Resistance and endurance remain active research frontiers.

\paragraph{Programmable logic in modern flows.} HDL synthesis maps \texttt{case} and \texttt{if}--\texttt{else} to MUXes and LUTs automatically;
the designer rarely instantiates PLA/PAL explicitly. Yet understanding them explains why a 4-LUT FPGA needs $\lceil n/4\rceil$ LUTs
per $n$-input function and why wide AND-OR structures benefit from dedicated carry/AND chains in modern FPGAs (e.g.\ Xilinx CARRY4).

\section{Exercises}
\label{sec:ch08-exercises}
\begin{enumerate}[label=E8.\arabic*, leftmargin=*, itemsep=0.8em]
    \item \textbf{SRAM vs.\ DRAM.} Compare 6T SRAM and 1T1C DRAM on transistor count, speed, refresh, and use case. Why is DRAM destructive on read?
    \item \textbf{Expansion.} Build $2\text{K}\times8$ from $1\text{K}\times8$ chips. How many chips? Draw address and data wiring, including $\overline{CS}$ decoding.
    \item \textbf{ROM logic.} Implement $F=\sum m(1,2,5,6)$ and $G=\sum m(0,3,7)$ (3 inputs) using an $8\times2$ ROM. List ROM contents. Repeat with a PLA using minimal product terms.
    \item \textbf{LUT mapping.} A 4-LUT can realise any 4-input function. How many 4-LUTs are needed for a 6-input AND? For a 4-bit adder's sum bits? Sketch the mapping.
    \item \textbf{Flash vs.\ EEPROM.} Distinguish NOR and NAND Flash on erase granularity, random access, and typical application (code storage vs.\ data storage).
\end{enumerate}
% End of Chapter 8

\backmatter
\end{document}
